% concatenated from arxiv_v1.tex, commands.tex, intro.tex, prelims.tex, main.tex, empirical.tex, related_work_conclusion.tex, proofs_and_more.tex
\documentclass[12pt]{article}
\usepackage[margin=1in]{geometry}
\usepackage[utf8]{inputenc}

\usepackage[english]{babel}
\usepackage[round]{natbib}
\bibliographystyle{apalike}

\usepackage[font=footnotesize]{caption}
\usepackage{mathtools,amssymb,amsthm}

\usepackage{subcaption}
\usepackage{pifont}
\usepackage{booktabs}

\usepackage{algorithm}
\usepackage{algorithmic}
\usepackage{wrapfig}
\usepackage{float}
\usepackage[toc,page]{appendix}


\usepackage[dvipsnames]{xcolor}
\definecolor{citeCol}{rgb}{0.5412, 0.2118, 0.0588}
\usepackage[colorlinks=true, citecolor=citeCol, linkcolor=citeCol, urlcolor=citeCol]{hyperref}

\usepackage[capitalize]{cleveref}
\usepackage{bm}
\usepackage{enumitem}

\usepackage{graphicx}
\usepackage{graphicx,amsfonts,mathrsfs}
\usepackage{amsthm,amsmath,amsfonts,amssymb,enumitem}
\usepackage{mathtools}
\theoremstyle{plain}
\newtheorem{ass}{Assumption}
\theoremstyle{plain}
\newtheorem{xmpl}{Example}
\usepackage{xcolor}
\usepackage{bm}
\usepackage{todonotes}


\newtheorem{lem}{Lemma}
\newtheorem{mydef}{Definition}
\newtheorem{prop}{Proposition}
\newtheorem{thm}{Theorem}
\newtheorem{cor}{Corollary}
\theoremstyle{remark}
\newtheorem{exa}{Example}


\newcommand{\bbR}{\mathbb{R}}
\newcommand{\R}{\mathbb{R}}
\newcommand{\Rz}{\mathbb{R}^{2d}}
\newcommand{\bbRM}{\mathbb{R}^m}
\newcommand{\bbRN}{\mathbb{R}^n}
\newcommand{\bbRMN}{\mathbb{R}^{m+n}}
\newcommand{\bbRP}{\mathbb{R}_{>0}}
\newcommand{\bbRD}{\mathbb{R}^d}
\newcommand{\N}{\mathbb{N}}
\newcommand{\Z}{\mathbb{Z}}

\newcommand{\cA}{\mathcal{A}}
\newcommand{\cB}{\mathcal{B}}
\newcommand{\cC}{\mathcal{C}}
\newcommand{\cF}{\mathcal{F}}
\newcommand{\cE}{\mathcal{E}}
\newcommand{\cD}{\mathcal{D}}
\newcommand{\cG}{\mathcal{G}}
\newcommand{\cH}{\mathcal{H}}
\newcommand{\cI}{\mathcal{I}}
\newcommand{\cJ}{\mathcal{J}}
\newcommand{\cP}{\mathcal{P}}
\newcommand{\cY}{\mathcal{Y}}
\newcommand{\cN}{\mathcal{N}}

\newcommand{\cX}{\mathcal{X}}
\newcommand{\cS}{\mathcal{S}}
\newcommand{\cT}{\mathcal{T}}
\newcommand{\cU}{\mathcal{U}}
\newcommand{\cV}{\mathcal{V}}
\newcommand{\cW}{\mathcal{W}}
\newcommand{\cL}{\mathcal{L}}

\newcommand{\tU}{\tilde{U}}
\newcommand{\tV}{\tilde{V}}
\newcommand{\tb}{\tilde{b}}
\newcommand{\tf}{\tilde{f}}
\newcommand{\tg}{\tilde{g}}
\newcommand{\tP}{\tilde{P}}
\newcommand{\tx}{\tilde{x}}
\newcommand{\ty}{\tilde{y}}
\newcommand{\tomega}{\tilde{\omega}}

\newcommand{\scrA}{\mathscr{A}}
\newcommand{\scrE}{\mathscr{E}}
\newcommand{\scrF}{\mathscr{F}}
\newcommand{\scrH}{\mathscr{H}}
\newcommand{\scrI}{\mathscr{I}}
\newcommand{\scrU}{\mathscr{U}}
\newcommand{\scrV}{\mathscr{V}}
\newcommand{\scrP}{\mathscr{P}}
\newcommand{\scrL}{\mathscr{L}}

\newcommand{\frkA}{\mathfrak{A}}

\newcommand{\bes}{\begin{equation*}}
\newcommand{\ees}{\end{equation*}}
\newcommand{\beas}{\begin{eqnarray*}}
\newcommand{\eeas}{\end{eqnarray*}}
\newcommand{\bea}{\begin{eqnarray}}
\newcommand{\eea}{\end{eqnarray}}
\newcommand{\be}{\begin{equation}}
\newcommand{\ee}{\end{equation}}
\newcommand{\bei}{\begin{itemize}}
\newcommand{\eei}{\end{itemize}}
\newcommand{\bec}{\begin{cases}}
\newcommand{\eec}{\end{cases}}
\newcommand{\ben}{\begin{enumerate}}
\newcommand{\een}{\end{enumerate}}


\newcommand{\Gm}{\Gamma}
\newcommand{\Om}{\Omega}
\newcommand{\veps}{\varepsilon}
\newcommand{\ve}{\varepsilon}
\newcommand{\vfi}{\varphi}

\newcommand{\bbD}{\mathbb{D}}
\newcommand{\bbP}{\mathbb{P}}
\newcommand{\bbE}{\mathbb{E}}
\newcommand{\bbF}{\mathbb{F}}
\newcommand{\bbQ}{\mathbb{Q}}
\newcommand{\bbS}{\mathbb{S}}
\newcommand{\bbW}{\mathbb{W}}
\newcommand{\bbX}{\mathbb{X}}
\newcommand{\bbV}{\mathbb{V}}


\newcommand{\bd}{\mathbf{d}}
\newcommand{\bD}{\mathbf{D}}
\newcommand{\bE}{\mathbf{E}}
\newcommand{\bg}{\mathbf{g}}
\newcommand{\bh}{\mathbf{h}}
\newcommand{\bt}{\mathbf{t}}
\newcommand{\bff}{\mathbf{f}}
\newcommand{\bp}{\mathbf{p}}
\newcommand{\bK}{\mathbf{K}}
\newcommand{\bT}{\mathbf{T}}
\newcommand{\bP}{\mathbf{P}}
\newcommand{\bQ}{\mathbf{Q}}
\newcommand{\bR}{\mathbf{R}}
\newcommand{\bw}{\mathbf{w}}
\newcommand{\bW}{\mathbf{W}}
\newcommand{\bY}{\mathbf{Y}}

\newcommand{\Int}{[0,1]}
\newcommand{\bbl}{\begin{block}}
\newcommand{\ebl}{\end{block}}
\newcommand{\ind}{\mathbf{1}}
\newcommand{\COND}{ \big \vert}
\newcommand{\dd}{\,\mathrm{d}}
\newcommand{\curv}{\kappa}

\newcommand{\hP}{\hat{P}}

\newcommand{\rmH}{\mathrm{H}}
\newcommand{\rmI}{\mathrm{I}}
\newcommand{\rmP}{\mathrm{P}}
\newcommand{\rmQ}{\mathrm{Q}}
\newcommand{\rmR}{\mathrm{R}}
\newcommand{\rme}{\mathrm{e}}

\newcommand{\bmX}{\bm{X}}
\newcommand{\bmY}{\bm{Y}}
\newcommand{\bmC}{\bm{C}}
\newcommand{\bmL}{\bm{L}}
\newcommand{\bmZ}{\bm{Z}}
\newcommand{\bmu}{\bm{u}}
\newcommand{\bmw}{\bm{w}}
\newcommand{\bmf}{\bm{f}}
\newcommand{\bmQ}{\bm{Q}}


\newcommand{\ip}[2]{\langle #1,#2\rangle}
\newcommand*{\defeq}{\coloneqq}

\newcommand{\msf}[1]{\mathsf{#1}}
\newcommand{\D}{\mathrm{d}}
\newcommand{\eps}{\varepsilon}




 \newcommand{\msa}{\mathsf{A}}
 \newcommand{\bb}{\mathbb}

 \newcommand{\rset}{\mathbb{R}}
 \newcommand{\bfc}{\mathbf{c}}
 \newcommand{\KLd}{\mathrm{KL}}
 \def\eqsp{\;}
 \def\rset{\mathbb{R}}


\def\ie{\textit{i.e.}}
\def\cadlag{càdlàg}
\def\eqsp{\;}
\newcommand{\coint}[1]{\left[#1\right)}
\newcommand{\ocint}[1]{\left(#1\right]}
\newcommand{\ooint}[1]{\left(#1\right)}
\newcommand{\ccint}[1]{\left[#1\right]}
\newcommand{\cointLigne}[1]{[#1)}
\newcommand{\ocintLigne}[1]{(#1]}
\newcommand{\oointLigne}[1]{(#1)}
\newcommand{\ccintLigne}[1]{[#1]}
\newcommand{\cointLine}[1]{[#1)}
\newcommand{\ocintLine}[1]{(#1]}
\newcommand{\oointLine}[1]{(#1)}
\newcommand{\ccintLine}[1]{[#1]}

\newcommand{\Idd}{\operatorname{I}_d}
\newcommand{\gauss}{\mathbf{N}}
\newcommand{\rmd}{\mathrm{d}}


\newcommand{\tcr}[1]{\textcolor{red}{#1}}

 \newcommand{\alain}[1]{\todo[color=blue!20]{{\bf AD:} #1}}


\newcommand{\PE}{\mathbb{E}}
\newcommand{\1}{\bm{1}}


\usepackage{pifont}

\newcommand{\cmark}{\ding{51}}
\newcommand{\xmark}{\ding{55}}
\newcommand{\gray}[1]{\textcolor{gray}{#1}}

\renewcommand{\thefootnote}{\fnsymbol{footnote}}

\definecolor{blue}{rgb}{0,.1, .6}
\definecolor{royalblue}{rgb}{0.2549, 0.4118, 0.8824}
\newcommand{\blue}[1]{\textcolor{royalblue}{#1}}
\newcommand{\aram}[1]{\textcolor{purple}{#1}}
\newcommand{\suggest}[1]{\royalblue{[Suggestion: #1]}}
\definecolor{stdgray}{gray}{0.48}
\providecommand{\wentry}[2]{#1\,{\color{stdgray}$\pm$\,#2}}

\theoremstyle{plain}
\newtheorem{theorem}{Theorem}[section]
\newtheorem{proposition}[theorem]{Proposition}
\newtheorem{lemma}[theorem]{Lemma}
\newtheorem{corollary}[theorem]{Corollary}
\theoremstyle{definition}
\newtheorem{definition}[theorem]{Definition}
\newtheorem{assumption}[theorem]{Assumption}
\theoremstyle{remark}
\newtheorem{remark}[theorem]{Remark}
\newtheorem{example}[theorem]{Example}



\title{Trajectory inference via Acceleration Matching}

\author{
\begin{tabular}{cc}
Bartolo Dazzini\thanks{Department of Mathematics, University of Padova. \texttt{bartolo.dazzini@studenti.unipd.it}}
&
Giovanni Conforti\thanks{Department of Mathematics, University of Padova. \texttt{giovanni.conforti@unipd.it}}
\\[0.6em]
Alain Durmus\thanks{CMAP, \'Ecole polytechnique. \texttt{alain.durmus@polytechnique.edu}}
&
Aram-Alexandre Pooladian
\thanks{Institute for Foundations of Data Science, Yale University. \texttt{aram-alexandre.pooladian@yale.edu}}
\thanks{Corresponding author}
\end{tabular}
}

\date{\today}

\begin{document}
\maketitle
\abstract{
Trajectory inference is a fundamental problem in many scientific domains: given a collection of unpaired snapshots of observations at discrete time points, the goal is to generate smooth trajectories that best resemble and interpolate the data.
Existing algorithms exhibit computational challenges:
they either rely on preprocessing subroutines to enforce smoothness or on simulation-based training objectives, both of which can be expensive.
In order to overcome these limitations, we propose a new algorithm called Acceleration Matching (\texttt{AM}). Our approach consists of lifting the original interpolation problem to phase space and then regressing onto an explicit conditional acceleration field that induces random, smooth trajectories that agree with the prescribed marginals. Importantly, our resulting training algorithm only requires positional data, avoids trajectory simulation during training, and is devoid of expensive preprocessing. We provide ample numerical evidence suggesting that \texttt{AM} is competitive with or superior to existing algorithms on several benchmark problems from the existing literature.
}


\section{Introduction}
We investigate the trajectory inference problem: given snapshots of unpaired observations collected at multiple time points
$
0=t_0<t_1<\cdots<t_J=1$, the goal is to generate trajectories that fit the given marginal data. In probabilistic terms, we want to find a measure-valued curve $t \in \ccint{0,1} \mapsto \mu_t \in {\cal P}(\R^d)$  such that
\begin{align}\label{eq:marg_constraints}
     \mu_{t_j} = \rho_{t_j} \quad \forall j \in \{0,1,\ldots,J\}\,,
\end{align}
for the given (empirical) probability distributions $\{\rho_{t_j}\}_{j=0}^J$.

This problem appears in a variety of data-driven research domains, including astronomy, economics, and single-cell RNA sequencing, among others \citep{franzke2015stochastic,kazakevivcius2021understanding,yeo2021generative,bunne2022proximal,kubica2007efficient,liounis2020autonomous}. In each of these cases, data are obtained over time in the form of unpaired observations. For example, the destructive nature of single-cell RNA sequencing makes it impossible to have repeated measurements of the same cell over time as the population (of cells) continues to evolve. As a result, the nature of the data renders standard procedures (based on, for example, time-series modeling) unhelpful. These unique challenges have led to several mathematical and algorithmic developments for the problem; see~\citet{lavenant2024toward, hong2025trajectory, bunne2023learning, schiebinger2019optimal}, among many others.

A naive approach is to simply construct $(\mu_t)_{t \in [0,1]}$ in a piecewise manner, learning $J$ independent segments of the probability path of interest. This can be accomplished by separately learning transport maps between each $\rho_{t_j}$ and $\rho_{t_{j+1}}$ and then stitching the paths together. For instance, one can learn optimal transport maps \citep{manole2021plugin,pooladian2021entropic,deb2021rates} or instead pass through the popular flow matching (FM) framework \citep{lipman2022flow,albergo2022building,liu2022flow}.
\begin{figure}[t]
    \centering
    \begin{subfigure}[t]{0.48\linewidth}
        \centering
        \includegraphics[width=\linewidth]{figures/piecewise_trajectories_fm.png}
    \end{subfigure}
    \hfill
    \begin{subfigure}[t]{0.48\linewidth}
        \centering
        \includegraphics[width=\linewidth]{figures/smooth_trajectories_am.png}
    \end{subfigure}
    \caption{\texttt{AM} outputs smoother trajectories \textbf{(right)} than piecewise flow matching \textbf{(left)}.}
    \label{fig:piecewise_v_splines}
\end{figure}

Depending on the method, the learn-then-stitch approach unfortunately requires, in the worst case, solving $J$ separate learning problems. Moreover, without additional preprocessing that explicitly enforces smooth dynamics \citep{rohbeck2025modeling}, this naive approach can produce trajectories that lack global smoothness, leading to non-smooth, potentially unrealistic trajectories. The left side of Figure~\ref{fig:piecewise_v_splines} displays trajectories given by training a simple vector field using the flow matching (FM) framework; the trajectories are visibly non-smooth due to the abrupt discontinuity in the velocity field as it switches between the marginals.

Our goal, by contrast, is to generate smooth positional trajectories by construction. To this end, we lift the original problem and consider trajectories given by stochastic differential equations (SDEs) on phase space (see Section~\ref{sec:phasespace} for more background), which are of the form
\begin{align}\label{sde_pp}
    \begin{cases}
        \dd X_t = V_t \dd t \\
        \dd V_t = a_t(X_t, V_t) \dd t + \sqrt{\varepsilon} \dd B_t\,,
    \end{cases}
\end{align}
where $a_t : \R^{2d}\to\R^d$ is an acceleration field and $(B_t)_{t\in[0,1]}$ is standard Brownian motion with diffusion parameter $\eps > 0$. These dynamics generate smooth trajectories by design, as the velocities now have continuous sample paths and thus the positions vary smoothly. Our approach then consists of defining and targeting an appropriate acceleration field that yields dynamics which satisfy the marginal constraints ${\rm Law}(X_{t_j}) = \rho_{t_j}$ for all $j \in \{0,1,\ldots,J\}$.

While we are not the first to consider dynamics of the form \eqref{sde_pp}, related prior works require simulating trajectories as part of their training pipeline and, as a result, are computationally expensive even for low-dimensional datasets \citep{chen2023deep,theodoropoulos2025momentum}.\footnote{Here we simply mean that trajectory rollouts are required somewhere in the training algorithm; we make no distinction between rollouts that require backpropagation and those that do not.} Summarizing our observations so far, it stands to reason that existing algorithms for trajectory inference come with multiple computational burdens. They either are not entirely simulation-free, do not explicitly enforce smooth dynamics or do so only through ad hoc preprocessing; see Section~\ref{sec:related_work} for a thorough comparison.

\paragraph{Contributions.}
In this work, we present \emph{Acceleration Matching} (\texttt{AM}), a novel approach for trajectory inference. Our algorithm is based on the Markovianization of kinetic Brownian bridges in phase space (see Lemma~\ref{lem:bridge_process} and Theorem~\ref{thm:main} in Section~\ref{sec:condacc} for more precise details) and can thus be viewed as an extension of the  flow matching framework \citep{lipman2022flow,albergo2022building,liu2022flow} for learning second-order dynamics in the presence of multiple marginal constraints. Upon constructing a suitable, entirely explicit acceleration field, we naturally obtain a training objective that only assumes access to positional data, is guaranteed to satisfy the constraints \eqref{eq:marg_constraints}, and most importantly, is entirely simulation-free.
From the right-hand side of Figure~\ref{fig:piecewise_v_splines}, we see that trajectories generated by \texttt{AM} are visibly smoother than those generated by piecewise FM.
We demonstrate the computational effectiveness of \texttt{AM} on several standard datasets from the literature in Section~\ref{sec:results}, where our approach is competitive or superior to existing algorithms that explicitly enforce smoothness constraints.




\section{Background}\label{sec:prelim}
\subsection{Notation}

\label{sec:notation}
Let $\cP(\R^d)$ denote the set of probability measures over $\R^d$. For a collection of probability distributions $\{P_i\}_{i=1}^M \subset \cP(\R^d)$, let $\Pi(P_1,\ldots,P_M)$ denote the set of couplings between the measures i.e., for $\pi \in \Pi(P_1,\ldots,P_M)$, $\pi \in \cP((\R^{d})^M)$ such that for every $i \in \{1,\ldots,M\}$, the $i^{\rm th}$ marginal of $\pi$ is $P_i$.
Let ${\cal C}^d \defeq C([0,1];\R^{d})$ be the space of continuous functions from the unit interval to $\R^{d}$ for $d \in \mathbb{N}$, endowed with the uniform topology and the corresponding Borel $\sigma$-field $\mathcal{B}({\cal C}^d)$. The set of probability measures over the resulting measurable space is denoted by $\mathcal{P}({\cal C}^d)$.
We use bold-font vectors (e.g., $\bm{t},\bm{x},\bm{v}$) to denote a concatenated vector; for example, $\bm{t} = (t_0,t_1,\ldots,t_J) \in \R^{J+1}$ and $\bm{x} = (x_0,x_1,\ldots,x_J) \in \R^{d(J+1)}$ (and similarly for $\bm v$).
For a set $A$, let $\bm{1}_A$ denote the $0$-$1$ indicator function over $A$.
For a smooth function $(x,v)\mapsto F(x,v)\in \R$, we write $\nabla_x F(x,v)$ and $\nabla_v F(x,v)$ to explicitly denote with respect to which argument the gradient is taken.





\subsection{Evolution of second-order probability flows}\label{sec:phasespace}
In this section, we briefly review probability paths on phase space, the joint space of positions $x\in\R^d$ and velocities $v \in \R^d$ of a given physical system. Throughout, we write phase space as $\rset^{2d}$.

Our approach is based on processes defined as solutions of second-order SDEs of the form
\begin{align}\label{eq:fp_dyn}
    \dd X_t^a = V_t^a \dd t\,, \quad \dd V_t^a = a_t(X_t^a,V_t^a)\dd t + \sqrt{\eps} \dd B_t\,, \quad t \in\ccint{0,1} \eqsp,
\end{align}
with initial conditions $(X_0^a,V_0^a)$,
where $(B_t)_{t \in [0,1]}$ is {$d$-dimensional} standard Brownian motion, $\eps > 0$, and $(t,x,v) \in \ccint{0,1}  \times \rset^{2d} \mapsto a_t(x,v) \in \R^d$. We then naturally interpret $(a_t)_{t\in[0,1]}$ as an inhomogeneous acceleration field that governs the underlying dynamics.


As an example, if $a_t(x,v) = - \gamma v$ for a parameter $\gamma > 0$, the resulting process is called \emph{kinetic underdamped Brownian motion}, denoted by $\msf R^{\gamma,\eps} \in \mathcal{P}({\cal C}^{2d})$. If $\gamma = 0$ (and thus $a_t \equiv 0$), then we simply refer to $\msf R^{0,\eps}$ as kinetic Brownian motion. Whenever there is no risk of confusion, we omit the explicit dependence on $\gamma \geq 0$ and $\eps > 0$ and simply write $\msf R$ for the reference process.

For any $t\in [0,1]$, let $\mu_t^a \defeq {\rm Law}(X_t^a,V_t^a)$ where $(X_t^a,V_t^a)_{t\in[0,1]}$ evolves according to \eqref{eq:fp_dyn}. Under appropriate conditions on $(a_t)_{t\in[0,1]}$ (see e.g., \cite{bogachev2022fokker}), it can be shown that the pair $(\mu_t^a,a_t)_{t\in[0,1]}$ satisfies the \emph{kinetic Fokker--Planck equation} (also called the Vlasov--Fokker--Planck equation), which is written as
\begin{align}\label{eq:kinetic_fp}\tag{$\msf{KFP}$}
\begin{split}
&\partial_t \mu_t^a + \langle \nabla_x \mu_t^a, v \rangle + \nabla_v \cdot(a_t \mu_t^a) = \frac{\eps}{2}\Delta_v \mu_t^a\,.
\end{split}
\end{align}



\section{Acceleration Matching}\label{sec:main}
In this section, we introduce Acceleration Matching (\texttt{AM}), our proposed algorithm for trajectory inference.

Our main technical observations and contributions appear in Section~\ref{sec:condacc}, where we construct a suitable interpolation process, the acceleration field corresponding to the second-order flow, and the Markovianized process. Our complete algorithm appears in Section~\ref{sec:loss}, which includes a detailed description of our simulation-free training objective. Finally, we describe the inference step of our method in Section~\ref{sec:gen_trajec}.

\subsection{Conditional acceleration fields and Markovianization}\label{sec:condacc}
Recall our second-order dynamics of interest given by \eqref{eq:kinetic_fp}.
Our goal is to define a particular family of acceleration fields that are easy to learn from positional data in the form of marginals $\{(t_j, \rho_{t_j})\}_{j=0}^J$. Importantly, we want the resulting dynamics to yield marginals that agree with \eqref{eq:marg_constraints}, i.e., $\mathrm{Law} (X_{t_j}^a) = \rho_{t_j}$.

To this end, take an arbitrary coupling $\pi \in \Pi(\rho_{t_0},\ldots,\rho_{t_J})$; for simplicity, consider the independent coupling of the marginals. Let $(X_t,V_t)_{t\in\ccint{0,1}} \sim \msf R$ be a kinetic (underdamped) Brownian motion for some fixed $\gamma \geq 0$ and $\varepsilon >0$. We define our interpolation process $\mathsf P^I \in \cP({\cal C}^{2d})$ in two steps:
\begin{enumerate}
    \item For $\bm{t} = (t_0,\ldots,t_J) \in [0,1]^{J+1}$, sample $X_{\bm t}  \sim \pi$
    \item Sample the remainder of the path from the regular conditional distribution $\mathsf R_{|\bm t}$,
\end{enumerate}
where a formal definition of $\msf R_{|\bm t}$ is given in Appendix~\ref{sec:supp_section_main}. Mathematically, our interpolation process is written as
\begin{align}\label{eq:interp-process}
    {\msf P}^I = \int \msf R_{|\bm t}(\bm x,\cdot)\,\pi(\rmd \bm x)\,,
\end{align}
where $\msf P^I \in {\cal P}({\cal C}^{2d})$. As we assume query access to $\pi$ (at least in the case of independent marginals), it remains to efficiently  sample from our conditional distribution $\msf R_{|\bm t}$. Thankfully, it turns out that this remaining step corresponds to sampling from a suitable Gaussian distribution. Indeed, our next result shows that if the initial velocity $V_0$ is Gaussian, then $V_{\bm t}$ given $X_{\bm t}$ is also Gaussian.
\begin{lem}
  \label{lem:condi_v_t}
    Assume that if $(X_0,V_0) \sim \msf R_0$, the conditional distribution of $V_0$ given $X_0$ is Gaussian almost surely with conditional mean $x_0 \mapsto m_0(x_0)$ and covariance matrix $x_0 \mapsto \Sigma_0(x_0)$. Then, the conditional distribution of $V_{\bm t}$ given $X_{\bm t}$ is also Gaussian with explicit conditional mean and covariance; see \eqref{eq:mean_cov_cond}.
\end{lem}
In brief, given only position knots $X_{\bm t}$, we can easily sample velocities $V_{\bm t}\,|\, X_{\bm t}$ in closed-form. The proof is an easy consequence of Lemma~\ref{lem:fund} and is deferred to Appendix~\ref{sec:proof-lemma-refl}.


The next step is to complete the sampling procedure for the \emph{entire path} $(X_t, V_t)_{t\in[0,1]}$ given the $J+1$ knots $(X_{\bm t}, V_{\bm t})$ in phase space. In the case of standard Brownian motion, this would be immediate, as it reduces to simulating independent Brownian bridges. Remarkably, an analogous characterization holds for the kinetic damped Brownian motion in phase space. For the remainder of this section, we focus on the case $\gamma = 0$, deferring the general case to the appendix.
\begin{lem}[Kinetic Brownian bridges]\label{lem:bridge_process}
Fix $t_j < t_{j+1} \in [0,1]$, $x_j,x_{j+1} \in \R^d$, $v_{j}, v_{j+1} \in \R^d$, and let $\msf R$ be kinetic Brownian motion with $\eps > 0$. {In the interval $[t_j,t_{j+1})$, the bridge of the reference process, i.e., the conditional distribution of $(X_t,V_t)_{t\in\ccint{t_j,t_{j+1}}}$ given  $(X_{t_{j}},V_{t_{j}}) = (x_j, v_j)$ and $(X_{t_{j+1}},V_{t_{j+1}}) = (x_{j+1}, v_{j+1})$} is the unique weak solution starting from $(X_{t_j},V_{t_j})$ to
\begin{align}\label{eq:bridge_proc}
        \dd X_t& = V_t\dd t \\
        \dd V_t &= a^j_t(X_t,V_t; X_{t_{j+1}},V_{t_{j+1}})\dd t + \sqrt{\varepsilon} \dd{ B_t^j}\,,
\end{align}
on $[t_j,t_{j+1})$, { where $\{(B_t^{k})_{t \in\ccint{t_k,t_{k+1}}} \, :\, k \in \{0,\ldots,J-1\}\}$ are independent $d$-dimensional Brownian motions}
and  for any $x,v,x_{j+1},v_{j+1} \in\rset^d$ and $t \in\ccint{t_j,t_{j+1}}$,
\begin{align}\label{eq:a_jt}
        a_t^{j}(x,v;x_{j+1},v_{j+1}) &=
    \frac{6(x_{j+1} - x)}{(t_{j+1}-t)^2} - \frac{2(v_{j+1} + 2v)}{(t_{j+1}-t)}\,.
\end{align}
\end{lem}
The proof of Lemma~\ref{lem:bridge_process} is based on an application of Doob's $h$-transform on phase space; see Appendix~\ref{sec:proof-bridge-process}. An important corollary of Lemma~\ref{lem:bridge_process} is that the non-Markovian path measure $\msf P^I$ can now be represented as a solution to a sequence of stochastic differential equations: given $(X_{\bm t}^I, V_{\bm t}^I)$, Lemma~\ref{lem:bridge_process} shows that for any  $j\in\{0,\ldots,J-1\}$, $(X^I_{t},V_t^I)_{t\in[t_j,t_{j+1}]}$ is the unique weak solution starting from $(X_{t_j},V_{t_j})$ to\looseness-1
\begin{align}\label{conditional_acceleration_SDE}
\begin{cases}
\dd X^I_t = V^I_t \dd t\\
\dd V^I_t = a_t^j(X^I_t,V^I_t ;X^I_{t_{j+1}},V^I_{t_{j+1}})\dd t + \sqrt{\eps}\dd B_t^j\,,
\end{cases}\raisetag{1.5\baselineskip}
\end{align}
over the interval $t\in[t_j,t_{j+1}]$.

Based on this observation, it becomes transparent that $(X_t^I,V_t^I)_{t\in\ccint{0,1}}$ is not Markov. However, we can find a diffusion process $(X_t^{M},V_{t}^M)_{t \in \ccint{0,1}}$ sharing the same marginals, replacing the acceleration field in \eqref{conditional_acceleration_SDE} with a suitable Markovianization that will be easy to learn using a parametric model. This is formalized by our next result.
\begin{thm}\label{thm:main}
  Consider $(X_t^I,V_t^I)_{t\in\ccint{0,1}} \sim \msf P^I$ be an interpolated process associated with a coupling $\pi \in \Pi(\rho_{t_0},\ldots,\rho_{t_J})$ and a kinetic Brownian distribution $\msf R$ for $\eps >0$ fixed.
Consider an acceleration field $(a_t)_{t\in\ccint{0,1}}$ satisfying almost surely
\begin{align*}
    a_t^{M}(X_{t}^I,V_t^I) \defeq \sum_{j=0}^{J-1} \bm{1}_{[t_j,t_{j+1})}(t) \, \bbE[a_t^j(X_t^I,V_t^I;X_{t_{j+1}}^I,V_{t_{j+1}}^I)|(X_t^I,V_t^I)]\,,
\end{align*}
where $a_t^j$ is given by \eqref{eq:a_jt}.  {In addition, suppose that $a_t^M$ is Lipschitz in all variables for any $t\in\ccint{0,1}$}. For $(X_0^M,V_0^M) \sim \mathsf P^I_0$, consider $(X_t^M,V_t^M)_{t\in\ccint{0,1}}$ the diffusion process solution over $ [0,1]$
\begin{align*}
    \begin{cases}
    \dd X_t^M = V_t^M \dd t \\
    \dd V_t^M = a_t^{M}(X_t^M,V_t^M)\dd t + \sqrt{\eps}\dd B_t\,.
    \end{cases}
\end{align*}
Then, for any $t \in [0,1]$, it holds that $\mathrm{Law}(X_t^M,V_t^M)=\mathrm{Law}(X_t^I,V_t^I)$, and therefore the positional marginals match the prescribed constraints.
\end{thm}
In brief, Theorem~\ref{thm:main} states that the marginal laws of the original non-Markovian interpolant can be recovered through a suitably defined acceleration field. This construction should be viewed as a direct \emph{lifted} analogue of, e.g., flow matching \citep{albergo2022building,lipman2022flow,liu2022flow} in the case $J=1$ (which corresponds to two marginals). Furthermore, in this way, we can impose multiple marginal constraints along the path, rather than only initial and terminal constraints.

\subsection{Loss function and training}\label{sec:loss}
We now describe a simple algorithm for learning the Markovian acceleration field from Theorem~\ref{thm:main}. Let $a_\theta$ denote a parametrized neural network such that $a_\theta : [0,1] \times \R^{2d} \to \R^d$. Given that $a_t^M$ is a (piecewise) conditional expectation, it is natural to consider the following population loss
\begin{align}\label{eq:pop_loss}
\begin{split}
        \cL_{\texttt{AM}}(a_\theta) = \sum_{j=0}^{J-1}\int_{t_j}^{t_{j+1}} \bbE\| a_\theta(t,X_t,V_t) - a_t^j(X_t,V_t;X_{t_{j+1}},V_{t_{j+1}})\|^2
        \dd t\,,
\end{split}
\end{align}
where the expectation is under $\mathsf P^I$.
It is therefore easy to see that $a_t^M$ is the unique minimizer of \eqref{eq:pop_loss}. Indeed, this holds by standard properties of conditional expectations as well as the piecewise nature of the acceleration field $a_t^M$ and so we omit the proof.

We now describe a practical implementation of \eqref{eq:pop_loss} in more detail; see also Algorithm~\ref{alg:stocacc}.
Recall that access to samples of positional data $X_{t_j}$ is readily available from the marginals $\{\rho_{t_j}\}_{j=0}^{J}$.
We then sample $V_{\bm t}\mid X_{\bm t}$ efficiently by Lemma~\ref{lem:condi_v_t} where we take ${\rm Law}(V_0\mid X_0) = \gauss(0,\sigma_v^2I)$ for $\sigma_v^2 >0$ for simplicity.
For each $j$, we first draw $s_j \sim {\rm Unif}((t_j, t_{j+1}))$, and again use Gaussian conditioning formulae to sample the kinetic Brownian bridge between consecutive knots; see Appendix~\ref{sec:bridge_sampling} for more details. {For uneven observation times, we set $\Delta_j = t_{j+1}-t_j$ and use it to weight the contribution from the $j^{\rm th}$ interval.} Note that the ``for loop'' in Algorithm~\ref{alg:stocacc} is merely for illustrative purposes; in practice this can be be parallelized or be replaced by sampling the indices.

We stress that unlike prior work, our loss does not merely decompose into a learning problem on disjoint intervals, as the velocities used to train $a_\theta$ are jointly sampled across time conditioned on all positional knots. This allows us to learn smooth dynamics purely from the formulation of the learning problem instead of explicitly imposing spline-like dynamics into the training pipeline.

\begin{algorithm}[t]
\caption{Training algorithm for Acceleration Matching}
\label{alg:stocacc}
\begin{algorithmic}[1]
\REQUIRE Marginals $\{\rho_{t_j}\}_{j=0}^J$; coupling $\pi \in \Pi(\rho_{t_0},\ldots,\rho_{t_J})$; network
$a_\theta : [0,1]\times \mathbb{R}^{2d} \to \mathbb{R}^d$
\WHILE{not converged}
    \STATE Sample $\bm{x} \sim \pi$
    \STATE Sample $\bm v \mid \bm x$ via Gaussian conditioning
    \FOR{$j \in \{0,1,\ldots,J-1\}$}
    \STATE {Set $\Delta_j = t_{j+1}-t_j$}
    \STATE Sample $s_j \sim \mathrm{Unif}((t_j,t_{j+1}))$
    \STATE Sample $(x_{s_j},v_{s_j}) \mid (\bm x_{t_j},\bm v_{t_j}),(\bm x_{t_{j+1}},\bm v_{t_{j+1}})$ via Gaussian conditioning
    \ENDFOR
    \STATE Compute loss $\widehat{\mathcal{L}}_{\mathtt{AM}}(a_\theta) = \sum_{j=0}^{J-1}{\Delta_j}\|a_\theta(s_j,x_{s_j},v_{s_j}) - a_{s_j}^j(x_{s_j},v_{s_j}; \bm x_{t_{j+1}}, \bm v_{t_{j+1}})\|^2$
    \STATE Update parameters $\theta \gets \texttt{optim}(\theta, \nabla_\theta
    \widehat{\mathcal{L}}_{\mathtt{AM}}(a_\theta))$
\ENDWHILE
\STATE \textbf{Return} $a_\theta$
\end{algorithmic}
\end{algorithm}








\subsection{Generating trajectories}\label{sec:gen_trajec}
We now describe how to generate trajectories given a learned acceleration field obtained following Algorithm~\ref{alg:stocacc}. The procedure has two stages: first, we learn an initial velocity sampler, then we integrate the learned dynamics in phase space.

Theorem~\ref{thm:main} requires the initial phase space pair to satisfy
\begin{align*}
      (X_0^M,V_0^M) \sim \msf P_0^I
\end{align*}
in order for the marginals generated by the Markovian acceleration field to agree at all times. At inference, however, only the initial position law $X_0\sim\rho_0$ is prescribed. In order to approximate the missing conditional law ${\rm Law}(V_0\mid X_0)$, we will train a conditional sampler $q_\phi(\cdot\mid x_0)$. We obtain valid training data $(X_0^{(i)},V_0^{(i)})_{i=1}^N$ by first sampling knots $(X_{t_0},\ldots,X_{t_J})$, then sampling the corresponding conditional velocity $V_0\,|\, X_{\bm t}$ as before, retaining only the initial position and velocity.


For simplicity, one can take $q_\phi$ to be a (conditional diagonal) Gaussian,
\[
    q_\phi(\cdot\mid x_0)
    =
    \gauss(m_\phi(x_0),\Sigma_\phi(x_0)),
\]
where $m_\phi$ and $\Sigma_\phi$ are outputs of a neural network $\phi$ that takes in $x_0$ as input, and we train the weights of $\phi$ (and thus $m_\phi$ and $\Sigma_\phi$) via maximum likelihood:
\begin{align*}
    \min_\phi \sum_{i=1}^N
    \frac{1}{2}
    \langle
    v_0^{(i)} - m_{\phi}(x_0^{(i)}),
    \Sigma_{\phi}(x_0^{(i)})^{-1}
    \bigl(v_0^{(i)} - m_{\phi}(x_0^{(i)})\bigr)
    \rangle
    +
    \frac{1}{2}\log\det \Sigma_{\phi}(x_0^{(i)})\,.
\end{align*}
See Appendix~\ref{sec:cond_sampling} for more details. While other conditional sampling methods can also be used (such as those proposed by \cite{baptista2024conditional}), we found this simple Gaussian model sufficient for our purposes.

Given a new $x_0\sim\rho_0$, we then sample $v_0\sim q_\phi(\cdot\mid x_0)$ and generate trajectories by discretizing the phase-space SDE in \eqref{sde_pp} using, for example, a standard Euler--Maruyama scheme; see Algorithm~\ref{alg:em_sampler}.

\begin{algorithm}[t]
\caption{Euler--Maruyama sampler}
\label{alg:em_sampler}
\resizebox{\linewidth}{!}{
\begin{minipage}{\linewidth}
\begin{algorithmic}[1]
\REQUIRE Initial state $x_0 \sim \rho_0$, learned conditional sampler $q_\phi(\cdot|x_0)$; learned acceleration field
$\widehat a_\theta$; noise scale $\varepsilon$; number of steps $K$
\STATE Draw $v_0 \sim q_\phi(\cdot\, |\, x_0)$ and set $h = 1/K$
\FOR{$k=0,\ldots,K-1$}
    \STATE Update position $ x_{(k+1)h}
        \gets
        x_{kh} + h\, v_{kh}$
    \STATE Draw $\xi_k \sim \mathcal{N}(0,I)$
    \STATE Update velocity $v_{(k+1)h}
        \gets
        v_{kh}
        + h \, \widehat a_\theta(kh, x_{kh}, v_{kh})
        + \sqrt{\varepsilon h}\,\xi_k$
\ENDFOR
\STATE \textbf{Return} $(x_{kh})_{k=0}^{K}$
\end{algorithmic}
\end{minipage}
}
\end{algorithm}


\section{Numerical experiments}\label{sec:results}
We now present a set of experiments comparing \texttt{AM} to other state-of-the-art methods for the task of trajectory inference. Our main comparisons are with 3MSBM \citep{theodoropoulos2025momentum} (which is also based on a lifting to phase space) and simulation-free flow-matching methods such as MMFM~\citep{rohbeck2025modeling}, OT-CFM~\citep{tong2023improving,pooladian2023multisample}, and OT-MFM~\citep{kapusniak2024metric}.

We consider three variants of MMFM; see Section~\ref{sec:related_work} for more details. The first version uses a constant noise coefficient, denoted $\text{MMFM}_{\rm const}$, which induces Gaussian smoothing at the marginal knots $\rho_{t_j}$ rather than matching the marginals exactly.
The second variant, denoted $\text{MMFM}_{\rm van}$, uses a time-varying noise coefficient that vanishes at the marginal knots and therefore interpolates the prescribed marginals.\footnote{Both of these variants appear in the official implementation by \cite{rohbeck2025modeling}. In particular, the majority of the experiments use $\text{MMFM}_{\rm const}$.} For $\text{MMFM}_{\bullet}$, the training parameters are all the same as our method. We also use the official \href{https://github.com/Genentech/MMFM}{MMFM repository}, with all default parameters, as a baseline for the low-dimensional experiments, which we denote simply by $\text{MMFM}$\@. We were {unable to reproduce the results} from the \href{https://github.com/panostheo98/3MSBM}{3MSBM repository} (nor had success implementing the method ourselves) and so we instead report values from their paper directly. For the remaining baselines, we will report the results from the corresponding papers, denoted by $\text{OT-CFM}^*$ and $\text{OT-MFM}^*$.

As in prior work, we consider datasets of the form
\begin{align*}
    \{(t_j,\widehat\rho_{t_j})\}_{j=0}^J,
\end{align*}
where each $\widehat\rho_{t_j}$ is an empirical measure consisting of $n_{t_j}$ atoms in $\R^d$. For each dataset, we train on a subset of the available marginals and evaluate whether the generated trajectories match held-out marginals. Following the conventions of prior work, we use the $2$-Wasserstein metric as the evaluation metric for the low-dimensional datasets and the $1$-Wasserstein metric for the high-dimensional single-cell datasets. Unless otherwise stated, we average results over five trials.

\subsection{Low-dimensional datasets}
We first evaluate \texttt{AM} on two low-dimensional datasets widely used in the literature: ocean-current trajectories in the Gulf of Mexico (GoM), and Lotka--Volterra (LV) dynamics~\citep{shen2024multi}. For the GoM and LV datasets, $J=8$ (thus nine total marginals) and we train on the even-time marginals, holding out the odd times $(t_1,t_3,t_5,t_7)$.
At inference time, we sample initial positions from $\widehat\rho_{t_0}$, simulate trajectories forward in time, and compare the generated samples with both training and held-out marginals using the $2$-Wasserstein distance. In Figures~\ref{fig:traj_gom} and~\ref{fig:traj_lv}, we denote the held-out marginals by $t_j^*$.





\begin{table}[t]
\centering
\caption{Average training and holdout $W_2$ errors for GoM and LV datasets; lower is better. }
\label{tab:low-dimensional-w2-300-official}
\begin{tabular}{lcccc}
& \multicolumn{2}{c}{GoM} & \multicolumn{2}{c}{LV} \\
\cmidrule(lr){2-3}\cmidrule(lr){4-5}
Method & Holdout $W_2$ & Train $W_2$ & Holdout $W_2$ & Train $W_2$ \\
\midrule
3MSBM$^{*}$ & 0.135 & --- & 0.230 & --- \\
MMFM & $0.209 \pm 0.012$ & $0.104 \pm 0.015$ & $0.427 \pm 0.031$ & $0.400 \pm 0.052$ \\
\midrule
MMFM$_{\rm const}$ & $0.191 \pm 0.018$ & $0.125 \pm 0.019$ & $0.530 \pm 0.036$ & $0.423 \pm 0.071$ \\
MMFM$_{\mathrm{van}}$ & $0.298 \pm 0.084$ & $0.131 \pm 0.048$ & $1.116 \pm 0.539$ & $0.809 \pm 0.467$ \\
\textbf{AM (Ours)} & $0.163 \pm 0.010$ & $0.093 \pm 0.022$ & $0.376 \pm 0.035$ & $0.246 \pm 0.029$ \\
\end{tabular}
\end{table}
Table~\ref{tab:low-dimensional-w2-300-official} summarizes our results for the low-dimensional datasets. \texttt{AM} achieves the lowest holdout and training error on both GoM and LV compared to all MMFM variants, though it falls short compared to the reported results from 3MSBM.

Figures~\ref{fig:traj_gom} and~\ref{fig:traj_lv} show representative simulated trajectories for GoM and LV\@. In both datasets, MMFM$_{\rm const}$ has trajectories that are heavily concentrated around the mean of each marginal. As desired, the trajectories generated by MMFM$_{\rm van}$ appear to match the training marginals but otherwise appear to be unstable between them. In contrast, \texttt{AM} produces trajectories that more consistently match the marginals along the entire path.

\begin{figure}[t!]
    \centering
    \includegraphics[width=0.85\linewidth]{figures/traj_gom.png}
    \caption{Visualization of trajectories generated for the GoM dataset.}
    \label{fig:traj_gom}
\end{figure}

\begin{figure}[h!]
    \centering
    \includegraphics[width=0.85\linewidth]{figures/traj_lv.png}
    \caption{Visualization of trajectories generated for the LV dataset.}
    \label{fig:traj_lv}
\end{figure}


\subsection{High-dimensional datasets}
We next evaluate \texttt{AM} on higher-dimensional datasets that arise from single-cell biology. Following prior work, we use PCA representations of developing embryoid bodies (EB)~\citep{moon2019eb} and CITE-seq data (CITE)~\citep{burkhardtcite}, reducing the dimensionality to either $5$ or $50$. As in prior work, we use a leave-one-marginal-out (LOO) protocol, always keeping the initial and final marginals in the training set, and report the average held-out $W_1$ error. For EB5, there are 5 marginals and we leave out one of $t_1, t_2,$ or $t_3$, and average the resulting errors; for CITE5 and CITE50, we leave out either $t_1$ or $t_2$, and average the resulting errors.

\begin{table}[t]
\centering
\caption{Average LOO $W_1$ error on EB5, CITE5, and CITE50; lower is better.}
\label{tab:summary-high-dim}
\begin{tabular}{@{}lccc@{}}
\textbf{Method} & \textbf{EB5} & \textbf{CITE5} & \textbf{CITE50} \\
\midrule
OT-CFM$^{*}$ & \wentry{0.790}{0.068} & \wentry{0.882}{0.058} & \wentry{38.756}{0.398} \\
OT-MFM$^{*}$ & \wentry{0.713}{0.039} & \wentry{0.724}{0.070} & \wentry{36.394}{1.886} \\
\midrule
MMFM$_{\mathrm{const}}$ & \wentry{0.824}{0.089} & \wentry{0.542}{0.018} & \wentry{41.674}{0.399} \\
MMFM$_{\mathrm{van}}$ & \wentry{0.971}{0.143} & \wentry{1.074}{0.077} & \wentry{200.946}{19.987} \\
{AM} & \wentry{0.791}{0.074} & \wentry{0.634}{0.021} & \wentry{49.831}{ 1.809} \\
\end{tabular}
\end{table}





The high-dimensional results are summarized in Table~\ref{tab:summary-high-dim}. On EB5, \texttt{AM} improves over both MMFM variants and is comparable to OT-CFM, while OT-MFM reports a slightly lower error. On CITE5, \texttt{AM} performs better than most of its competitors, though ultimately MMFM$_{\rm const}$ returns the lowest error. The CITE50 setting is the most challenging: while \texttt{AM} is far more stable than MMFM$_{\rm van}$ and is somewhat competitive with MMFM$_{\rm const}$, all three of our implementations appear to lag behind the reported errors from the other flow-matching approaches. We compare generated trajectories of \texttt{AM} and MMFM in Appendices~\ref{app:eb} and~\ref{app:cite}.

\begin{figure}[t]
    \centering
    \includegraphics[width=0.95\linewidth]{figures/eb5_l1o_t1_pca.png}
    \caption{Visualization of generated marginals by \texttt{AM} and MMFM$_{\rm van}$ and MMFM$_{\rm const}$ on the EB5 dataset in the LOO setting ($t_1$ is left out).}
    \label{fig:eb5-t1-pca}
\end{figure}




\subsection{Related work}\label{sec:related_work}

\subsubsection*{Multi-sample flow matching and related algorithms}
Flow matching \citep{lipman2022flow,albergo2022building,liu2022flow} is a simulation-free approach for generating probability paths between two marginals.
Multi-marginal flow matching (MMFM) is a simulation-free approach proposed by \citet{rohbeck2025modeling}, which uses ideas reminiscent of other works that incorporate data-dependent couplings  \citep{tong2023improving,pooladian2023multisample,albergo2023stochastic} combined with using cubic splines.
Similar to the proposal by \citet{chewi2021fast}, they first compute pairwise optimal transport matchings between the marginals. For datasets of moderate size, these matchings can be obtained using a standard optimal transport package \citep{flamary2021pot} where, writing $n = \max_j n_{t_j}$ as the maximum number of points per marginal, the overall preprocessing cost is ${\cal O}(Jn^3)$ and the storage cost is therefore ${\cal O}(Jn^2)$. (Note that this preprocessing step can be expensive for $n \gtrsim 10^4$.) Then, given such a pairwise aligned coupling, they sample data $X_{0:J} \defeq (X_0,\ldots,X_J) \sim {\rm pairwise}(\rho_0,\ldots,\rho_J)$, and approximate the following velocity field\looseness-1
\begin{align*}
    \bbE[\partial_t \,{\rm spline}(t; X_{0:J}) + \sigma_t' \xi\,|\, {\rm spline}(t;X_{0:J}) + \sigma_t \xi = z]\,,
\end{align*}
where ${\rm spline}(t; X_{0:J})$ is a cubic spline between the drawn knots, $\sigma_t \geq 0$ is the noise coefficient, and $\xi \sim \cN(0,I)$. Upon learning the resulting vector field, they simulate the resulting ordinary differential equation to obtain the trajectories. \citet{rohbeck2025modeling} were also interested in learning conditional velocity fields via classifier-free guidance, which we do not pursue here.

The code from \citet{rohbeck2025modeling} mostly uses a constant variance, i.e., $\sigma_t' = 0$, this is what we call ${\rm MMFM}_{\rm const}$. They also consider $\sigma_t = 4(t_{j+1}-t)(t - t_j)/(t_{j+1}-t_j)^2$ which vanishes at the marginal endpoints; we denote this by ${\rm MMFM}_{\rm van}$. As we mentioned before, the former induces trajectories where the positional marginals are Gaussian-smoothed versions of the true marginal knots, whereas the latter matches the marginals exactly.

Optimal transport conditional flow matching (OT-CFM) uses minibatch optimal transport couplings between the marginals when fitting a velocity field \citep{tong2023improving,pooladian2023multisample}. \cite{kapusniak2024metric} proposed OT-Metric Flow Matching (OT-MFM), which incorporates a Riemannian metric from the data in addition to learning a vector field. The Riemannian metrics take the form of kernel density estimators (e.g., radial basis functions).

\subsubsection*{Multi-marginal Schr\"odinger bridge methods}
 Many neural-network-based algorithms for solving the trajectory inference problem are based on multi-marginal Schr\"odinger bridges. For instance, multi-marginal stochastic flow matching (MMSFM) \citep{lee2025multi} is a variant of MMFM that incorporates some ideas from \citet{shi2024diffusion}, though still uses learning pairwise computed OT couplings (sometimes triplets). Schr\"odinger bridge with iterative reference refinement (SBIRR) \citep{shen2024multi} requires some degree of trajectory simulation as part of the algorithm and has non-smooth positional trajectories as they use standard SDEs, not lifted ones, for generating trajectories.

The work of \citet{theodoropoulos2025momentum} also passes to phase space (which they refer to as ``momentum''), where their goal is to learn a Markovian acceleration field. Apart from this similarity, our proposed method is quite distinct from theirs. They regress on a different conditional acceleration field, one which depends on \emph{several} future positions in a non-trivial way. Their proposed conditional acceleration field is defined as
\begin{align*}
    &a_t^{\rm{3MSBM}}((x_t,v_t)\,|\,\{\bm x_{j+1} \, : \, t_{j+1} \geq t\}) = C_1^j(t)(x_t - \bm x_{j+1}) + C_2^j(t) v_t + C_3^j(t)\sum_{\ell=j+1}^J \lambda_\ell \bm x_\ell\,,
\end{align*}
for $t \in [t_j, t_{j+1})$, where $\{\bm x_j\}_{j=0}^J$ denotes the sampled knots, $\{\{C_k^j(t)\}_{k=1}^3\}_{j=0}^J$ are time-varying coefficients corresponding to each segment, and $\{\lambda_j\}_{j=0}^J$ are also computed. Note that for any choice of $J$, the time-varying coefficients need to be computed manually, and are not as explicit as ours (recall \eqref{eq:a_jt}). Additionally, their algorithm ultimately requires simulating trajectories (forward and backward in time) using their conditional acceleration field for each update step. While they do not require function calls to their parametric function approximator, the simulation of SDEs per iteration during training is nevertheless computationally undesirable, and ultimately not simulation-free.

\section{Conclusion}
We propose Acceleration Matching as a principled, scalable algorithm for trajectory inference. To learn smooth trajectories, we lift the learning problem to phase space, with the goal of learning an appropriate Markovianized acceleration field. This perspective yields a fully explicit, regression objective for learning dynamics with multiple marginal constraints that, unlike existing algorithms, requires minimal preprocessing, incurs little per-iteration overhead, and is entirely simulation-free during training. Several natural directions remain beyond the scope of this work, such as investigating the role of the reference process, improving overall algorithmic complexity, extending our framework to instance-specific Riemannian manifolds~\citep{chen2023flow,yim2023se,haviv2024wasserstein}, quantifying discretization guarantees in the spirit of~\citet{chen2022sampling} and~\citet{conforti2025kl}, and applying our method to more computationally intensive settings where either $J$ or $d$ (or both) are large.


\section*{Acknowledgments}
AAP is grateful for financial support from the Institute for Foundations of Data Science at Yale University and to Ricardo Baptista for helpful comments on an earlier version of this draft.


\bibliography{main}

\newpage
\appendix
\section{Supplement of Section~\ref{sec:main}}
\label{sec:supp_section_main}
We denote by $(X_t)_{t\in[0,1]}$ the canonical process, i.e., $X_t(h) = h_t \in \R^d$ for any $h \in {\cal C}^d$ and any $t \in [0, 1]$. For any $t$, $\msf Q_t$ denotes the $t$-marginal $(X_t)_\sharp \msf Q$ for $\msf Q \in \mathcal{P}({\cal C}^d)$. By a slight abuse of notation, we denote by $(X_t,V_t)_{t\in[0,1]}$ the canonical process on the phase space $\rset^{2d}$, i.e., $X_t(h) = \mathrm{proj}_x(h_t)$ and $V_t(h) = \mathrm{proj}_v(h_t)$ for any $h \in {\cal C}^{2d}$ and any $t \in [0, 1]$. Here, $\mathrm{proj}_x$ and $\mathrm{proj}_v$ denote the projections onto the first and last $d$ components, respectively.

We also recall the dynamics associated with kinetic underdamped Brownian motion, denoted $\msf R = \msf R^{\gamma,\eps}$
\begin{align}\label{eq:ref_proc}
    \begin{cases}
        \dd X_t = V_t\dd t \\
        \dd V_t = -\gamma V_t \dd t + \sqrt{\varepsilon} \dd B_t\,.
    \end{cases}
\end{align}
For additional shorthand, let $Z_t \defeq (X_t,V_t)$ for $t \in [0,1]$.

We now provide a formal definition of $\msf R_{|\bm t}$; this is a Markov kernel on $(\rset^{d})^{J+1} \times \mathcal{B}(\mathcal{C}^{2d})$ satisfying for any measurable sets $A \in \mathcal{B}(\mathcal{C}^{2d})$ and $B \in \mathcal{B}((\rset^d)^{J+1})$, $\msf R(A\cap \{X_{\bm t} \in B\}) = \PE[\1_{B}(X_{\bm t}) \msf R_{|\bm t}(X_{\bm t},A)]$.



\subsection{Proof of Lemma~\ref{lem:condi_v_t}}
\label{sec:proof-lemma-refl}
\begin{lemma}\label{lem:fund}
  Conditionally on $(X_0,V_0)$, the reference process \eqref{eq:ref_proc} is a Gaussian process
  with mean and covariance functions given by, for $s,t \in \ccint{0,1}$ and $s \leq t$,
  \begin{equation*}
    m_Z(t) = \begin{pmatrix}
X_0 + V_0  a_t \\[6pt]
V_0 b_t
\end{pmatrix}  =  A_t Z_0 \,,
    \end{equation*}
where $a_t = \gamma^{-1}(1-\rme^{-\gamma t})$, $b_t = \rme^{-\gamma t}$, and
\begin{align*}
    {A}_t &= \begin{pmatrix}
        1 & a_t \\
        0 & b_t
      \end{pmatrix} \otimes \Idd \,.
\end{align*}
Moreover, it holds that for $s\leq t$, $\Gamma_Z(s,t) \defeq {\rm Cov}(Z_s,Z_t\,|\,Z_0)$ is given by
\begin{align*}
     \Gamma_Z(s,t) &=
     \begin{pmatrix}
\mathrm{Var}(X_s) + a_{t-s}\mathrm{Cov}(X_s,V_s)  & b_{t-s} \mathrm{Cov}(X_s,V_s)  \\
\mathrm{Cov}(X_s,V_s) + a_{t-s}  \mathrm{Var}(V_s)   & b_{t-s}  \mathrm{Var}(V_s)
     \end{pmatrix}
\otimes \Idd \,,
\end{align*}
where
\begin{align*}
    \mathrm{Var}(V_s) &= (\eps /(2 \gamma))( 1  - \rme^{-2 \gamma s} )\,, \\
    \mathrm{Var}(X_s)& =   (\varepsilon/\gamma^2) \left[ s - 2 a_s + (1/(2\gamma))(1-\rme^{-2\gamma s})\right]\,, \\
  \mathrm{Cov}(X_s,V_s) &  =  (\eps/2)a_s^2\,.
\end{align*}
In particular, if $\gamma = 0$, then the above simplifies to
  \begin{align*}
    m_Z(t) = \begin{pmatrix}
        X_0 + V_0 t \\
        V_0
    \end{pmatrix}
    = A_t Z_0, \qquad A_t =
    \begin{pmatrix}
        1 & t \\
        0 & 1
    \end{pmatrix} \otimes \Idd,
\end{align*}
and
\begin{align*}
    \Gamma_Z(s,t) =
\varepsilon     \begin{pmatrix}
  s^2(3t-s)/6 & s^2/2\\
  s^2/2 +s(t-s) & s
    \end{pmatrix}
    \otimes \Idd.
\end{align*}
\end{lemma}

\begin{proof}
First, we only show the result for $d=1$ since given $(X_0,V_0)$ all components are independent and therefore the result tensorizes. Note that for $s,t \in \ccint{0,1}$, we can readily solve the SDE for the velocity term to yield
\begin{align*}
    V_t = \rme^{-\gamma(t-s)}V_s + \eta_{s,t}^V \defeq \rme^{-\gamma(t-s)}V_s + \sqrt{\eps}\int_s^t \rme^{-\gamma(t-u)}\dd B_u\,,
\end{align*}
and for $s=0$, we have that $V_t = \rme^{-\gamma t}V_0 + \sqrt{\eps}\int_0^t \rme^{-\gamma(t-u)} \dd B_u$. Solving for the position variable then yields
\begin{align*}
    X_t = X_s + \int_s^t V_u \dd u  &= X_s + a_{t-s}V_s + \eta_{s,t}^X\\
    &\defeq X_s + a_{t-s}V_s+ \sqrt{\varepsilon} \int_s^t \int_{s}^{u} \rme^{-\gamma(u - r)} \dd B_{r} \dd u\,,
\end{align*}
where in the second step we used that $\int_s^t \rme^{-\gamma(u-s)}\dd u = a_{t-s}$. Therefore, as a solution of a linear SDE, $(X_t,V_t)_{t\in \ccint{0,1}}$ is indeed a Gaussian process given initial datum $(X_0,V_0)$. We now proceed and compute the remaining explicit expressions, which make heavy use of the It\^o isometry and the Markov property of the Brownian motion.

To start, we compute the double integral as
\begin{align*}
  \int_s^t \int_s^u \rme^{-\gamma(u-r)} \dd B_r \dd u = \int_s^t a_{t-r}\dd B_r\,,
\end{align*}
which further simplifies the expression for $X_t$, which will be useful later.

We readily compute the means of the Gaussian process as
\begin{align*}
    \bbE[V_t\,|\,(X_0,V_0)] = \rme^{-\gamma t} V_0\,, \quad \bbE[X_t\,|\,(X_0,V_0)] = X_0 + a_tV_0\,,
\end{align*}
which implies the representation $m_Z(t) = A_t Z_0$ and the accompanying matrix $A_t$.

Conditionally on $(X_0,V_0)$, define the centered variables
\begin{align*}
    \overline{V}_s = \sqrt{\eps}\int_0^s \rme^{-\gamma(s-u)}\dd B_u\,, \quad \overline{X}_s = \sqrt{\eps}\int_0^s a_{s-u}\dd B_u\,,
\end{align*}
and then we readily compute
\begin{align*}
    &{\rm Var}(V_s) = \bb E\Bigl[\overline{V}_s^2\Bigr] = \eps \int_0^s \rme^{-2\gamma(s-u)}\dd u = \frac{\eps}{2\gamma}(1-\rme^{-2\gamma s})\,,
\end{align*}
and similarly
\begin{align*}
{\rm Var}(X_s) = \bb E\Bigl[\overline{X}_s^2\Bigr] = \eps \int_0^s a_{s-u}^2 \dd u  &= \frac{\eps}{\gamma^2}\int_0^s (1 - 2\rme^{-\gamma r} + \rme^{-2\gamma r}) \dd r \\
&= \frac{\eps}{\gamma^2}\bigl(s - 2a_s + (2\gamma)^{-1}(1-\rme^{-2\gamma s}) \bigr)\,,
\end{align*}
and also
\begin{align*}
    {\rm Cov}(X_s,V_s) = \eps \int_0^s \rme^{-\gamma(s-u)} a_{s-u} \dd u = \frac{\eps}{2}a_s^2\,.
\end{align*}
We now compute the covariance between times $s \leq t$. Note that $V_t$ can be written in terms of $V_s$ and an additional Brownian motion term that is independent of $V_s$ (and similarly for $X_s$). Due to this independence, we can compute the covariance between times as
\begin{align*}
    {\rm Cov}(V_s,V_t) = {\rm Cov}(V_s, \rme^{-\gamma (t-s)}V_s + \eta_{s,t}^V) = {\rm Cov}(V_s, \rme^{-\gamma (t-s)}V_s) = \rme^{-\gamma (t-s)}{\rm Var}(V_s)\,,
\end{align*}
and similarly
\begin{align*}
    & {\rm Cov}(X_s,V_t) = {\rm Cov}(X_s, \rme^{-\gamma (t-s)}V_s) = \rme^{-\gamma(t-s)}{\rm Cov}(X_s,V_s)\,, \\
    &{\rm Cov}(V_s,X_t) = {\rm Cov}(V_s, X_s + a_{t-s}V_s) = {\rm Cov}(X_s,V_s) + a_{t-s}{\rm Var}(V_s)\,, \\
    &{\rm Cov}(X_s,X_t) = {\rm Cov}(X_s, X_s + a_{t-s}V_s) = {\rm Var}(X_s) + a_{t-s}{\rm Cov}(X_s,V_s)\,,
\end{align*}
which concludes the proof.
\end{proof}

Recall our notation conventions $0=t_0<t_1<\cdots<t_J= 1$ with $\bm{t} = (t_0,\ldots,t_J)$. Additionally define $\bm a, \bm b \in \R^{J+1}$ as
\begin{align*}
    \bm a \defeq a_{\bm t} \defeq (a_{t_0}, \ldots, a_{t_J}) \,, \quad
    \bm b \defeq b_{\bm t} \defeq (\rme^{-\gamma t_0}, \ldots, \rme^{-\gamma t_J})\,,
\end{align*}
as well as $\bm q \defeq (\bm a,\bm b) \in \R^{2(J+1)}$, and recall 
\begin{align*}
V_{\bm{t}}
=
\begin{pmatrix}
V_0&
V_{t_1}&
\hdots&
V_{t_J}
\end{pmatrix}^{\top}
\in\mathbb R^{(J+1)d}\eqsp,
\qquad
X_{\bm{t}}
=
\begin{pmatrix}
X_{0}&
X_{t_1}&
\hdots&
X_{t_J}
\end{pmatrix}^{\top}
\in\mathbb R^{(J+1)d}\eqsp.
\end{align*}

\begin{lemma}\label{lem:condi_v_t_detailed}
Let $X_0 = x_0 \in \R^d$ be fixed, and assume that
\begin{align*}
    V_0 \mid X_0 = x_0 \sim {\gauss}(m_0(x_0), \Sigma_0(x_0))\,.
\end{align*}
Then conditionally on $X_0 = x_0$, $(X_{\bm t}, V_{\bm t})$ is (jointly) Gaussian with mean $\bm m = (m_{\bm X}, m_{\bm V})$, where the components are
\begin{align}
    m_{\bm X} \defeq \bm 1_{J+1} \otimes x_0 + \bm a \otimes m_0(x_0)\,,\quad m_{\bm V} = \bm b \otimes m_0(x_0)\,,
\end{align}
and covariance $\bm{\Sigma} \in \R^{2(J+1)d  \times 2(J+1)d}$ given by
\begin{align*}
    \bm{\Sigma} \defeq
    \begin{pmatrix}
\Sigma_{\bm X} &\Sigma_{\bm{XV}}\\
\Sigma_{\bm{VX}}&\Sigma_{\bm{V}}
\end{pmatrix} \defeq
({\bm{q}} {\bm{q}}^{\top} \otimes \Sigma_0(x_0))  + \begin{pmatrix}
K_{\bm X} &K_{\bm{VX}}^\top\\
K_{\bm{VX}} &K_{\bm V}
\end{pmatrix} \otimes I_d\,,
\end{align*}
where for $i,k \in \{0,\ldots,J\}$ we define
\begin{align*}
  [K_{\bm X}]_{ik} & = \mathrm{Var}(X_{t_i\wedge t_k}) + a_{|t_i-t_k|} \mathrm{Cov}(X_{t_i\wedge t_k}, V_{t_i\wedge t_k})\,, \\
[K_{\bm V}]_{ik}
  &= \rme^{-\gamma |t_i-t_k|} \mathrm{Var}(V_{t_i\wedge t_k})\,, \\
[K_{\bm{VX}}]_{ik}
  &    =              \begin{cases}
                    \mathrm{Cov}(X_{t_i},V_{t_i})
+
a_{t_k-t_i}\mathrm{Var}(V_{t_i}) & i \leq k\,, \\
\rme^{-\gamma(t_i-t_k)} \mathrm{Cov}(X_{t_k},V_{t_k}) & i >k \,.
                  \end{cases}
\end{align*}
Moreover, $V_{\bm t} \mid X_{\bm t}$ is Gaussian with mean and covariance given by
\begin{align}\label{eq:mean_cov_cond}
\begin{split}
m_{\bm V\mid \bm X} &= {m}_{\bm V} +
\Sigma_{\bm{VX}}\Sigma_{\bm X}^{\dagger}
(X_{\bm{t}}- m_{\bm X})\,, \\
\Sigma_{\bm V\mid \bm X}
&=
\Sigma_{\bm{V}}
-
\Sigma_{\bm{VX}}\Sigma_{\bm{X}}^\dagger\Sigma_{\bm{XV}}\,,
\end{split}
\end{align}
where $\Sigma_{\bm{X}}^{\dagger}$ denotes the Moore--Penrose pseudoinverse of $\Sigma_{\bm{X}}$.
\end{lemma}
\begin{proof}
By Lemma~\ref{lem:fund}, given $(X_0,V_0)$, the vector $( X_{\bm{t}},V_{\bm{t}})$ is jointly
Gaussian, so all that remains is to compute the mean and covariance.

Conditional on $X_0 = x_0$, we have by the representations in Lemma~\ref{lem:fund} that
\begin{align}\label{eq:lemma_eq}
    X_{t_i} = x_0 + a_{t_i} V_0 + \eta_i^X\,, \quad V_{t_i} = b_{t_i} V_0 + \eta_i^V\,,
\end{align}
where we recall that we defined $b_t = \rme^{-\gamma t}$. We further write the random variable $V_0\mid X_0 = x_0$ as $V_0 = m_0(x_0) + \xi_0$ for $\xi_0 \sim {\gauss}(0,\Sigma_0(x_0))$. Taking expectations leads to the definitions of $m_{\bm X}$ and $m_{\bm V}$ since $\xi_0$, $\eta_i^X$, and $\eta_i^V$ have mean zero. Then rewriting \eqref{eq:lemma_eq}, we have
\begin{align*}
    (X_{\bm t}, V_{\bm t}) = (m_{\bm X}, m_{\bm V}) + (\bm q \otimes I_d)\xi_0 + (\eta_{\bm t}^X, \eta_{\bm t}^V)\,,
\end{align*}
where, to compute the covariance, we separate the two sources of noise (which are independent from one another). We directly compute that
\begin{align*}
    {\rm Cov}((\bm q \otimes I_d)\xi_0) = (\bm q \otimes I_d){\rm Cov}(\xi_0)(\bm q \otimes I_d)^\top = (\bm q \bm q^\top)\otimes \Sigma_0(x_0)\,,
\end{align*}
and for ${\rm Cov}(\eta_{\bm t}^X, \eta_{\bm t}^V)$, we appeal to the existing formulae in Lemma~\ref{lem:fund}. The remaining conditional Gaussian formulas follow from standard computations that we omit.
\end{proof}


We end this section by computing the pseudoinverse of $\Sigma_{\bm{X}}$.
\begin{lemma}
  \label{lem:pseudo_inverse}
Let $\gamma \geq 0, \eps > 0$ and consider the reference process $\mathsf R = \mathsf R^{\gamma,\eps}$, and recall the notation and setting of Lemma~\ref{lem:condi_v_t_detailed}. Write $L_{\bm X} \in \R^{J\times J}$ such that $[L_{\bm X}]_{ik} = [K_{\bm X}]_{ik}$ for $i,k \in \{1,\ldots,J\}$. Then $L_{\bm X}$ is positive definite, and it holds that
\begin{align}\label{eq:cov_pinv}
    (\Sigma_{\bm X})^\dagger = \begin{pmatrix}
        0 & 0 \\
        0 & \Theta^{-1}
    \end{pmatrix}\,,
\end{align}
where, for $\tilde {\bm a} = (a_{t_1}, \ldots, a_{t_J}) \in \R^J$, $\Theta$ is defined as
\begin{align*}
    \Theta = \tilde{\bm a} \tilde{\bm a}^\top \otimes \Sigma_0(x_0) + L_{\bm X}\otimes I_d\,,
\end{align*}
and its inverse is given by
\begin{align}\label{eq:theta_t_inv}
    \Theta^{-1} = L_{\bm X}^{-1} \otimes I_d - (L_{\bm X}^{-1} \tilde{\bm a}\tilde{\bm a}^\top L_{\bm X}^{-1})\otimes [\Sigma_0(x_0)(I_d + c_\star \Sigma_0(x_0))^{-1}]
\end{align}
for  $c_\star = \langle \tilde{\bm a}, L_{\bm X}^{-1} \tilde{\bm a}\rangle$. In particular, if $\Sigma_0(x_0)$ is invertible, then
\begin{align*}
    \Theta^{-1} = L_{\bm X}^{-1} \otimes I_d - (L_{\bm X}^{-1} \tilde{\bm a}\tilde{\bm a}^\top L_{\bm X}^{-1})\otimes (\Sigma_0(x_0)^{-1} + c_\star I_d)^{-1}\,.
\end{align*}
\end{lemma}
\begin{proof}
We first prove that $L_{\bm X}$ is positive definite. As before, work over $\R$ since the $d$-dimensional covariance tensorizes with $I_d$.

From the representation in \eqref{eq:lemma_eq}, the centered positional noise at time $t_i$ is
\begin{align*}
    \eta_i^X = \sqrt{\eps}\int_0^{t_i} a_{t_i-u}\dd B_u\,,
\end{align*}
where $a_r = \gamma^{-1}(1-\rme^{-\gamma r})$ for $\gamma>0$, with the convention $a_r=r$ when $\gamma=0$. Hence, for any $c=(c_1,\ldots,c_J)\in\R^J$,
\begin{align*}
        \langle c, L_{\bm X} c\rangle
    =
    \eps\int_0^{1}
    \left(
        \sum_{i=1}^J c_i a_{t_i-u}\bm 1_{\{u\leq t_i\}}
    \right)^2
    \dd u \geq 0\,.
\end{align*}
We proceed by contradiction: if $\langle c, L_{\bm X} c\rangle = 0$ for $c \neq 0$, then the integrand would be zero for almost every $u$. Let $m$ be the largest index such that $c_m\neq 0$; on the interval $(t_{m-1},t_m)$, all terms with $i<m$ vanish because $u>t_i$, and all terms with $i>m$ have zero coefficient by the definition of $m$. Thus the integrand reduces to $(c_m a_{t_m-u})^2$, which is strictly positive. This yields a contradiction, so $L_{\bm X}$ is positive definite and thus invertible.

The block form of $(\Sigma_{\bm X})^\dagger$ follows because $t_0=0$ makes the first row and column of $\Sigma_{\bm X}$ vanish, while the remaining block is $\Theta$. Since $L_{\bm X}$ is positive definite and $\Sigma_0(x_0)$ is positive semidefinite, $\Theta$ is invertible. The displayed expression for $\Theta^{-1}$ is then the Sherman--Morrison--Woodbury formula applied to $\Theta=L_{\bm X}\otimes I_d + \tilde{\bm a}\tilde{\bm a}^\top\otimes \Sigma_0(x_0)$.
\end{proof}

\subsection{Proof of Lemma~\ref{lem:bridge_process}}\label{sec:proof-bridge-process}
\begin{lemma}
For $\gamma \geq 0, \eps > 0$, consider  kinetic underdamped Brownian motion $\mathsf R = \mathsf R^{\gamma,\eps}$, and for $t_j, t_{j+1} \in [0,1]$ with $t_j < t_{j+1}$, fix endpoints $(x_j,v_j)$ and $(x_{j+1},v_{j+1})$ respectively. Then the bridge of the reference process conditioned on $(X_{t_j},V_{t_j}) = (x_j,v_j)$ and $(X_{t_{j+1}},V_{t_{j+1}}) = (x_{j+1},v_{j+1})$ is the unique solution to
\begin{align}\label{eq:bridge_sde_general}
    \begin{cases}
        \dd X_t = V_t\dd t \\
        \dd V_t = (\eps\nabla_v \log p_{t_{j+1}-t}(X_t,V_t,x_{j+1},v_{j+1}) - \gamma V_t)\dd t + \sqrt{\varepsilon} \dd B_t\,,
    \end{cases}
\end{align}
on $t \in [t_j,t_{j+1})$, where $p_{t_{j+1} - t}(x,v;x_{j+1},v_{j+1})$ is the transition density of the reference process.
\end{lemma}

\begin{proof}
Since both the $\msf R$-bridge and the solution of \eqref{eq:bridge_sde_general} are Markov processes, it suffices to show that their infinitesimal generators coincide. If we denote by $p_t((x,v),(x',v'))$ the transition density of the reference process $\msf R$, then the marginal flow of the bridge at time $t \in [t_j, t_{j+1})$ is
\begin{equation*}
\mu^{\msf br}_t(x,v) \propto
{p_{t-t_j}((x_j,v_j),(x,v))p_{t_{j+1}-t}((x,v),(x_{j+1},v_{j+1}))}\,.
\end{equation*}
Our goal is to understand the evolution of $\mu^{\msf{br}}_t$ in the weak sense. To this end, consider now a smooth test function $\phi:[t_{j},t_{j+1}]\times \mathbb{R}^{2d}\to\mathbb{R}$ with compact support; we compute
\begin{align*}
    \partial_t \int \phi_t \mu_t^{\msf{br}} = \int (\partial_t \phi_t) \mu_t^{\msf{br}} + \phi_t (\partial_t p_{t-t_j})p_{t_{j+1}-t} + \phi_t p_{t-t_j}(\partial_t p_{t_{j+1}-t})\,.
\end{align*}
To continue, we require the infinitesimal generator of the kinetic (underdamped) reference process, which we recall is given by
\begin{align*}
    \cL_\gamma f(x,v) = \langle v, \nabla_x f(x,v)\rangle - \gamma \langle v, \nabla_v f(x,v)\rangle + \frac{\eps}{2}\Delta_{v}f(x,v)
\end{align*}
for a test function $f$. Now, we use the appropriate forward and backward Kolmogorov equations to see that
\begin{align*}
\partial_t p_{t-t_j} = \cL_\gamma^* p_{t-t_j}\,, \quad \text{and} \quad \partial_t p_{t_{j+1}-t} = -\cL_\gamma p_{t_{j+1}-t}
\end{align*}
resulting in
\begin{align}
\begin{split}\label{eq_h_trans_1}
    \partial_t \int \phi_t \mu_t^{\msf{br}} &= \int (\partial_t \phi_t) \mu_t^{\msf{br}} + \phi_t ( \cL_\gamma^* p_{t-t_j})p_{t_{j+1}-t}  - \phi_t p_{t-t_j}(\cL_\gamma p_{t_{j+1}-t}) \\
    &= \int (\partial_t \phi_t) \mu_t^{\msf{br}} + \cL_\gamma (\phi_t p_{t_{j+1}-t}) p_{t-t_j}  - \phi_t p_{t-t_j}(\cL_\gamma p_{t_{j+1}-t})\,,
\end{split}
\end{align}
where the last line uses integration by parts with the adjoint $\cL_\gamma$.
Next, we use the commutation identity for two test functions \citep{bakry2014analysis}
\begin{equation*}
\begin{aligned}
\mathcal{L}_\gamma(g h) &=(\mathcal{L}_\gamma g) h  + g(\mathcal{L}_\gamma h) + \varepsilon \langle \nabla_{v} g, \nabla_v h\rangle
\end{aligned}
\end{equation*}
to rewrite \eqref{eq_h_trans_1} as
\begin{align}\label{eq_h_trans_2}
\begin{split}
\partial_t \int \phi_t \mu_t^{\msf{br}} &=
\int \Bigl((\partial_t \phi_t)\mu_t^{\mathsf{br}} +  (\mathcal{L}_\gamma \phi_t) p_{t_{j+1}-t} p_{t-t_{j}} + \varepsilon\, \langle \nabla_v \phi_t, \nabla_v p_{t_{j+1}-t} \rangle p_{t - t_j}\Bigr) \\
 &= \int ( (\partial_t + \mathcal{L}_\gamma) \phi_t   + \varepsilon \, \langle \nabla_v \log p_{t_{j+1}-t}, \nabla_v \phi_t\rangle  )\mu_t^{\msf{br}} \,,
\end{split}
\end{align}
which implies that $\mu_t^{\msf{br}}$ is a weak solution to the Fokker--Planck equation associated with the generator
\begin{equation*}
\mathcal{L}^{\msf{br}}_{\gamma, t} := \cL_\gamma + \varepsilon  \langle \nabla_v \log p_{t_{j+1}-t}, \nabla_v \cdot \rangle\,.
\end{equation*}
\end{proof}

\begin{lemma}
Let $\mathsf R^{\gamma,\eps}$ be the kinetic underdamped Brownian motion reference process for $\gamma > 0$ and $\eps > 0$, and fix $T \in (0,1]$ and $z_T = (x_T,v_T) \in \R^{2d}$. Letting  $p_{T-t}((x,v);z_T)$ denote the transition density from time $t \in [0,T)$ to $T$ and $r \defeq T - t$, it holds that
\begin{align}\label{eq:damped_bridge_score}
    \eps \nabla_v \log p_{T-t}((x,v); z_T) =
    C_x(r;\gamma)(x_T - x - a_r v) + C_v(r;\gamma)(v_{T}- b_r v)\,,
\end{align}
where
\begin{align*}
    C_x(r;\gamma) &= \gamma^2(1-\rme^{-\gamma r})/D_r(\gamma)\,, \\
    C_v(r;\gamma) &= \gamma(\rme^{-2\gamma r} + 2 \gamma r \rme^{-\gamma r} - 1)/((1-\rme^{-\gamma r})D_r(\gamma))\,, \\
    D_r(\gamma) &= \gamma r(1 + \rme^{-\gamma r}) + 2\rme^{-\gamma r} - 2\,.
\end{align*}
In particular, $\gamma \to 0$ recovers \eqref{eq:a_jt}.
\end{lemma}
\begin{proof}
Writing $r = T-t$, it holds by Lemma~\ref{lem:fund} that, conditioning on $Z_t = (x,v)$
\begin{align*}
    (X_T,V_T) | (X_t,V_t) = (x,v) \sim {\gauss}\Bigl(\begin{pmatrix}
        x + a_rv \\
        b_r v
    \end{pmatrix},  Q_r\Bigr)\,,
\end{align*}
where $Q_r = \Gamma_Z(r,r)$ (and recall the notation $a_r = \gamma^{-1}(1-\rme^{-\gamma r})$ and $b_r = \rme^{-\gamma r}$ from Lemma~\ref{lem:fund}).
Thus up to a constant, we have that
\begin{align*}
&\log p_{T-t}((x,v),(x_T,v_T)) =\\
&\qquad -\frac{1}{2} \begin{pmatrix}
    x_T -
    (x + a_r v)\\
    v_T - b_r v
\end{pmatrix}^\top (Q_r \otimes I_d)^{-1}\begin{pmatrix}
    x_T -
    (x + a_r v)\\
    v_T - b_r v
\end{pmatrix}  + {\rm const}\,.
\end{align*}
As the matrix $Q_r^{-1}$ is explicit (via block 2-by-2 inversion), one can verify through a tedious calculation that indeed \eqref{eq:damped_bridge_score} holds, and setting $\gamma = 0$ recovers \eqref{eq:a_jt}.
\end{proof}
\subsection{Proof of Theorem~\ref{thm:main}}\label{proof_main}
Our proof follows the standard argument that passes through weak solutions of the kinetic underdamped Fokker--Planck equation. Letting $f : \R^d \times \R^d \to \R$ be a smooth test function, our goal is to compute
\begin{align*}
    \int f \,\partial_t \mu_t^I = \partial_t \Bigl(\int f \dd \mu_t^{I}\Bigr) = \partial_t \bbE[f(X_t^I,V_t^I)]\,,
\end{align*}
where $\mu_t^I$ is the density of $\mathsf P^I_t$. To proceed, we apply It\^o's lemma to first arrive at
\begin{align*}
        \dd f &= \langle \nabla_x f, \dd X_t^I \rangle + \langle \nabla_v f, \dd V_t^I\rangle + \frac{1}{2}\langle {\rm d} V_t^I, \nabla_{v}^2 f \dd V_t^I \rangle \\
        &= \bigl(\langle \nabla_x f, V_t^I \rangle + \langle \nabla_v f, a_t^j \rangle
        - \gamma \langle\nabla_v f, V_t^I \rangle +  \frac{\eps}{2}\Delta_v f\bigr)\dd t  + \sqrt\eps \langle \nabla_v f, \dd B_t\rangle\,.
    \end{align*}
Taking expectations, as the last term has mean zero by the martingale property, yields
\begin{align*}
    \partial_t \bbE[f(X_t^I,V_t^I)] &= \bbE\Bigl[ \langle \nabla_x f(X_t^I,V_t^I), V_t^I \rangle + \langle \nabla_v f(X_t^I,V_t^I), a_t^j(X_t^I,V_t^I) \rangle\\
    &\phantom{=}\quad - \gamma \langle\nabla_vf(X_t^I, V_t^I), V_t^I\rangle  + \frac{\eps}{2}\Delta_vf(X_t^I, V_t^I)\Bigr]\,,
\end{align*}
where recall that $a_t^j(X_t^I,V_t^I) = a_t^j(X_t^I, V_t^I; X_{t_{j+1}}^I, V_{t_{j+1}}^I)$. By the tower property of conditional expectations, it holds that
\begin{align*}
    &\bbE[\langle a_t^j(X_t^I, V_t^I; X_{t_{j+1}}^I, V_{t_{j+1}}^I), \nabla_v f(X_t^I,V_t^I) \rangle ] \\
    &\quad = \bbE[\bbE[\langle a_t^j(X_t^I, V_t^I; X_{t_{j+1}}^I, V_{t_{j+1}}^I), \nabla_v f(X_t^I,V_t^I) \rangle \,|\,(X_t^I,V_t^I)]] \\
    &\quad = \bbE[\langle \bbE[a_t^j(X_t^I, V_t^I; X_{t_{j+1}}^I, V_{t_{j+1}}^I)\,|\,(X_t^I,V_t^I)], \nabla_v f(X_t^I,V_t^I) \rangle ] \\
    &\quad = \bbE[\langle {a_t^{M}}(X_t^I,V_t^I), \nabla_v f(X_t^I,V_t^I) \rangle ]\,,
\end{align*}
where we recall that for an interval $[t_j,t_{j+1})$
\begin{align*}
    (t,x,v)\mapsto a_t^{M}(x,v) \defeq \bbE[a_t^j(X_t^I, V_t^I;X_{t_{j+1}}^I, V_{t_{j+1}}^I)\,|\,(X_t^I,V_t^I) = (x,v)]\,.
\end{align*}
Therefore, we see that the weak solution to $\mu_t^I$ satisfies
\begin{align*}
    \partial_t \bbE[f(X_t^I,V_t^I)] &= \bbE\Bigl[ \langle \nabla_x f(X_t^I,V_t^I), V_t^I \rangle + \langle \nabla_v f(X_t^I,V_t^I), a_t^{M}(X_t^I,V_t^I)\rangle \\
    &\phantom{=}\quad -\gamma \langle \nabla_v f(X_t^I, V_t^I), V_t^I \rangle + \frac{\eps}{2}\Delta_vf(X_t^I, V_t^I)\Bigr]\,.
\end{align*}
Since $\mu_0^M = \mu_0^I$ and, by assumption, the acceleration field implies a unique solution to the kinetic Fokker--Planck equation, we conclude the proof.


\section{Algorithmic subroutine details}
\subsection{Learning a conditional sampler}\label{sec:cond_sampling}
Recall that Theorem~\ref{thm:main} requires that $(X_0^M,V_0^M) \sim \msf P_0^I$, meaning that the initial position and velocity need to be appropriately sampled from the law of the interpolant at the initial time. Dropping the $M$ superscript, this (informally) amounts to sampling from
\begin{align*}
    p(x_0,v_0) = p(x_0) p(v_0\mid x_0)\,.
\end{align*}
We have sample access to $p(x_0)$ (this is the first marginal constraint), but we do not have access to the conditional law of $V_0$ given $X_0$; we remedy this through a simple variational approximation step.

Recall that we can efficiently sample interpolants
\begin{align*}
    v_{0:J}\mid X_{0:J} = x_{0:J}
\end{align*}
from the Gaussian conditioning formula \eqref{eq:mean_cov_cond}. We then extract the initial position and velocity which correspond to index $j=0$.
Repeating this procedure $n$ times returns a paired dataset $\{(x_0^{(i)}, v_0^{(i)})\}_{i=1}^n$. As a practical surrogate, we turn to maximum likelihood estimation with a Gaussian approximation. Letting $x\mapsto \phi(x)$ be a neural network that outputs a mean and diagonal covariance, i.e., $m_\phi(x)\in \R^d$ and $\Sigma_\phi(x) = {\rm diag}(s_\phi(x))$ for $s_\phi(x) \in \R^d_+$, we then optimize
\begin{align*}
    \max_\phi \sum_{i=1}^n \log q_\phi(v_0^{(i)}\mid x_0^{(i)})\,,
\end{align*}
or equivalently
\begin{align}
    \min_\phi \sum_{i=1}^n
    \frac{1}{2}\log\det \Sigma_\phi(x_0^{(i)})
    + \frac{1}{2}\langle v_0^{(i)} - m_\phi(x_0^{(i)}),
    \Sigma_{\phi}(x_0^{(i)})^{-1}\bigl(v_0^{(i)} - m_\phi(x_0^{(i)})\bigr)\rangle\,.
\end{align}

Table~\ref{tab:params-conditional-sampler} contains the hyperparameters used for training the neural network $\phi$, which has a simple MLP structure.

\subsection{Bridge sampling}\label{sec:bridge_sampling}
For a given $j \in \{0,1,\ldots,J-1\}$, define ${\cal I}_j = [t_{j},t_{j+1}]$ and $\delta_{j} = t_{j+1} - t_j$. We want to sample (a single) intermediate point (position and velocity) along the trajectory given two endpoint constraints
\begin{align*}
    (X_t,V_t) | (X_{t_{j+1}},V_{t_{j+1}}) = (x_{j+1},v_{j+1}), (X_{t_{j}},V_{t_{j}}) = (x_j,v_j)\,,
\end{align*}
where $t \in [t_j + c, t_{j+1} - c]$ for $0 < c \ll \delta_j/2$. Again, we can appeal to explicit formulae that arise in the study of Gaussian processes. For simplicity, we restrict our attention to the case $\gamma=0$.

Define
\begin{align*}
     Y_t = \begin{pmatrix}
        X_{t_j + t} \\
        V_{t_j + t}
    \end{pmatrix}\,,
\end{align*}
for $t \in [0,\delta_{j}]$, where the dynamics will again follow the reference process $\msf R$. Thus
\begin{align*}
     Y_t = m_t + \sqrt \varepsilon\, \xi_t = \begin{pmatrix}
        x_j + t v_j \\
        v_j
    \end{pmatrix} + \sqrt \varepsilon \begin{pmatrix}
        \int_0^t B_u \dd u \\
        B_t
    \end{pmatrix}\,.
\end{align*}
Here, the initialization pins down the left endpoint constraints and it remains to enforce the right endpoint. Since $(Y_t; Y_{\delta_{j}})$ is jointly Gaussian, we can readily compute ${\rm Cov}(\xi_t)$, ${\rm Cov}(\xi_{\delta_j})$ and ${\rm Cov}(\xi_t,\xi_{\delta_j})$ explicitly. By Gaussian conditioning, we have that $Y_t \, | \, Y_{\delta_j} \sim {\gauss}(m_{t|j}, C_{t|j})$ with
\begin{align*}
      m_{t|j} &= m_t + {\rm Cov}(\xi_t,\xi_{\delta_j}){\rm Cov}(\xi_{\delta_j})^{-1}(Y_{\delta_j} - m_{\delta_j})\,, \\
      C_{t|j} &= \varepsilon {\rm Cov}(\xi_t) - \varepsilon {\rm Cov}(\xi_t,\xi_{\delta_j}){\rm Cov}(\xi_{\delta_j})^{-1}{\rm Cov}(\xi_t,\xi_{\delta_j})^\top\,.
\end{align*}
To conclude, we simply substitute $Y_{\delta_j} = (x_{j+1},v_{j+1})$ in the above expression, and draw a sample. Note that this procedure can be efficiently parallelized.

\section{Remaining experimental details and results}
\subsection{Gulf of Mexico}\label{app:gom}
The Gulf of Mexico (GoM) dataset describes the movement of water particles in ocean currents. In our experiments, we work with preprocessed data from \citet{shen2024multi}. They used data from a HYbrid Coordinate Ocean Model (HYCOM) reanalysis \citep{hycom2024} at a specific time snapshot to compute the velocity field of the ocean currents. Then they generated $999$ trajectories in $\R^2$ and randomly sampled each of them at one of the nine time points. With this procedure, each particle is observed only at a single time point and each generated marginal contains $n=111$ samples. The dataset is divided into $5$ training marginals (the even times $t_0, t_2, t_4, t_6, t_8$) with the remaining $4$ times $(t_1, t_3, t_5, t_7)$ used for held-out evaluation. The data are centered and scaled to have zero mean and unit variance.


In Tables~\ref{tab:gom-training-marginals} and~\ref{tab:gom-eval-marginals}, we present the $2$-Wasserstein distance on training and held-out marginals, to have a more precise indicator of the performance of the algorithms relative to Table~\ref{tab:low-dimensional-w2-300-official}. As for the other datasets, precise training information can be found in Table~\ref{tab:params-acceleration}.

\begin{table}[t]
\centering
\caption{Per-timepoint $W_2$ at training marginals on the Gulf of Mexico dataset.}
\label{tab:gom-training-marginals}
\begin{tabular}{lcccc}
Method & $t_2$ & $t_4$ & $t_6$ & $t_8$ \\
\midrule
MMFM & \wentry{0.050}{0.008} & \wentry{0.089}{0.017} & \wentry{0.140}{0.028} & \wentry{0.240}{0.042} \\
MMFM$_{\mathrm{const}}$ & \wentry{0.071}{0.018} & \wentry{0.126}{0.040} & \wentry{0.167}{0.026} & \wentry{0.262}{0.028} \\
MMFM$_{\mathrm{van}}$ & \wentry{0.070}{0.025} & \wentry{0.127}{0.044} & \wentry{0.165}{0.083} & \wentry{0.294}{0.116} \\
\midrule
AM (Ours) & \wentry{0.094}{0.011} & \wentry{0.108}{0.046} & \wentry{0.125}{0.038} & \wentry{0.140}{0.036} \\
\end{tabular}
\end{table}

\begin{table}[t]
\centering
\caption{Per-timepoint $W_2$ at held-out marginals on the Gulf of Mexico dataset.}
\label{tab:gom-eval-marginals}
\begin{tabular}{lcccc}
Method & $t_1$ & $t_3$ & $t_5$ & $t_7$ \\
\midrule
3MSBM$^*$ & 0.20 & 0.18 & 0.07 & 0.09 \\
MMFM & \wentry{0.332}{0.010} & \wentry{0.187}{0.037} & \wentry{0.132}{0.030} & \wentry{0.187}{0.051} \\
\midrule
MMFM$_{\mathrm{const}}$ & \wentry{0.198}{0.021} & \wentry{0.192}{0.027} & \wentry{0.164}{0.050} & \wentry{0.210}{0.030} \\
MMFM$_{\mathrm{van}}$ & \wentry{0.285}{0.034} & \wentry{0.292}{0.125} & \wentry{0.295}{0.135} & \wentry{0.319}{0.172} \\
\midrule
AM (Ours) & \wentry{0.194}{0.009} & \wentry{0.182}{0.023} & \wentry{0.145}{0.067} & \wentry{0.130}{0.031} \\
\end{tabular}
\end{table}




\subsection{Lotka--Volterra}\label{app:lv}
This dataset, which simulates the Lotka--Volterra equations of a predator-prey model \citep{goel1971volterra}, is obtained from \citet{shen2024multi}. The data live in $\R^2$ and consist of $9$ marginals with $50$ points each. As with the GoM dataset, the even times are used for training ($t_0, t_2, t_4, t_6, t_8$) and the remaining $4$ for evaluation ($t_1, t_3, t_5, t_7$). Again, the data are centered and normalized. Tables~\ref{tab:lv-training-marginals} and~\ref{tab:lv-eval-marginals} respectively contain the training and evaluation errors in $W_2$ across the times of interest.

\begin{table}[t]
\centering
\caption{Per-timepoint $W_2$ at training marginals on the Lotka--Volterra dataset.}
\label{tab:lv-training-marginals}
\begin{tabular}{lcccc}
Method & $t_2$ & $t_4$ & $t_6$ & $t_8$ \\
\midrule
MMFM & \wentry{0.179}{0.031} & \wentry{0.223}{0.039} & \wentry{0.448}{0.051} & \wentry{1.150}{0.216} \\
MMFM$_{\mathrm{const}}$ & \wentry{0.168}{0.024} & \wentry{0.208}{0.018} & \wentry{0.464}{0.069} & \wentry{1.274}{0.296} \\
MMFM$_{\mathrm{van}}$ & \wentry{0.301}{0.082} & \wentry{0.509}{0.265} & \wentry{0.991}{0.713} & \wentry{2.244}{1.346} \\
AM (Ours) & \wentry{0.140}{0.042} & \wentry{0.193}{0.036} & \wentry{0.340}{0.144} & \wentry{0.559}{0.096} \\
\end{tabular}
\end{table}

\begin{table}[t]
\centering
\caption{Per-timepoint $W_2$ at held-out marginals on the Lotka--Volterra dataset.}
\label{tab:lv-eval-marginals}
\begin{tabular}{lcccc}
Method & $t_1$ & $t_3$ & $t_5$ & $t_7$ \\
\midrule
3MSBM$^*$ & 0.24 & 0.18 & 0.07 & 0.36 \\
\midrule
MMFM & \wentry{0.206}{0.039} & \wentry{0.211}{0.060} & \wentry{0.351}{0.017} & \wentry{0.941}{0.109} \\
MMFM$_{\mathrm{const}}$ & \wentry{0.534}{0.029} & \wentry{0.237}{0.034} & \wentry{0.364}{0.028} & \wentry{0.983}{0.134} \\
MMFM$_{\mathrm{van}}$ & \wentry{0.533}{0.266} & \wentry{0.861}{0.488} & \wentry{1.145}{0.697} & \wentry{1.927}{1.201} \\
\midrule
AM (Ours) & \wentry{0.520}{0.020} & \wentry{0.236}{0.074} & \wentry{0.238}{0.108} & \wentry{0.510}{0.115} \\
\end{tabular}
\end{table}



























\subsection{Embryoid body}\label{app:eb}
We use the preprocessed data from \cite{moon2019eb}, which are formed from the 100-dimensional PCA representation of the scRNA-seq of human embryoid cells collected across multiple time points. The dataset contains five marginals with 2,381, 4,163, 3,278, 3,665, and 3,332 cells (i.e., unpaired data) per marginal, respectively.

For our experiments, we use a lower-dimensional representation that considers only the first $5$ principal components (EB-5). Following the protocol from \cite{tong2020trajectorynet}, we perform a leave-one-out (LOO) procedure by training on four out of five marginals, normalizing the data to zero mean and unit variance. The evaluation is then performed at the held-out time (one of $t_1$, $t_2$, $t_3$) using the $1$-Wasserstein distance on normalized data. During training, we consider the full marginals and sample minibatches from them. At inference, we consider $2000$ points from the first marginal as starting points for the trajectories. Figure~\ref{fig:eb5-t1-pca} shows PCA projections of the ground truth marginals alongside those generated by each method. Table~\ref{tab:eb5-holdout} reports the ${W}_1$ distances on the held-out marginals under the leave-one-out protocol.

\begin{table}[t]
\centering
\caption{Mean and standard deviation of the $W_1$ error at held-out marginals on the EB-5 dataset.}
\label{tab:eb5-holdout}
\begin{tabular}{@{}lccc@{}}
Method & $t_1$ & $t_2$ & $t_3$ \\
\midrule
MMFM$_{\mathrm{const}}$ & \wentry{0.717}{0.026} & \wentry{0.862}{0.028} & \wentry{0.891}{0.068} \\
MMFM$_{\mathrm{van}}$   & \wentry{0.899}{0.055} & \wentry{0.890}{0.080} & \wentry{1.124}{0.128} \\
\midrule
AM & \wentry{0.763}{0.031} & \wentry{0.778}{0.015} & \wentry{0.831}{0.112} \\
\end{tabular}
\end{table}



\subsection{CITE}\label{app:cite}
This dataset \citep{burkhardtcite} consists of single-cell CITE-seq measurements collected at four time points over ten days. We considered both the $5$-dimensional (CITE5) and $50$-dimensional (CITE50) PCA projections, preprocessed following prior work by \cite{wang2026joint}. The training data are normalized to zero mean and divided by the maximum per-feature standard deviation, following \cite{neklyudov2024computational}. We evaluate all methods (considering the normalized space for CITE5 and the original PCA space for CITE50) under the leave-one-out (LOO) protocol as presented in \cite{kapusniak2024metric}, reporting $W_1$ on held-out marginals in Tables~\ref{tab:cite5-holdout} and~\ref{tab:cite50-holdout}. Figure~\ref{fig:cite5-t1-pca} shows the first PCA components of the ground truth and generated marginals using \texttt{AM} and MMFM.\looseness-1

\begin{figure}[h]
    \centering
    \includegraphics[width=0.8\linewidth]{figures/cite5_l1o_pca_new.png}
    \caption{Visualization of marginals generated by \texttt{AM}, MMFM$_{\rm van}$, and MMFM$_{\rm const}$ for the first PCA components on the CITE5 dataset in the LOO setting ($t_1$ is left out).}
    \label{fig:cite5-t1-pca}
\end{figure}



\begin{table}[t]
\centering
\caption{Per-timepoint $W_1$ at held-out marginals on the CITE5 dataset (LOO protocol).}
\label{tab:cite5-holdout}
\begin{tabular}{lcc}
Method & $t_1$ & $t_2$ \\
\midrule
MMFM$_{\mathrm{const}}$ & $0.530 \pm 0.003$ & $0.553 \pm 0.018$ \\
MMFM$_{\mathrm{van}}$   & $1.008 \pm 0.035$ & $1.140 \pm 0.044$ \\
AM & $0.633 \pm 0.011$ & $0.635 \pm 0.028$ \\
\end{tabular}
\end{table}

\begin{table}[t]
\centering
\caption{Per-timepoint $W_1$ at held-out marginals on the CITE50 dataset (LOO protocol).}
\label{tab:cite50-holdout}
\begin{tabular}{lcc}
Method & $t_1$ & $t_2$ \\
\midrule
MMFM$_{\mathrm{const}}$ & $41.665 \pm 0.228$ & $41.684 \pm 0.515$ \\
MMFM$_{\mathrm{van}}$   & $183.490 \pm 9.200$ & $218.402 \pm 10.238$ \\
AM & $48.093 \pm 0.318$ & $51.568 \pm 0.637$ \\
\end{tabular}
\end{table}









\clearpage
\newpage
\subsection{Further training details and hyperparameters}\label{app:training_details}
This section presents the parameters used for training the models, to allow for experimental reproducibility.

All models are trained with the Adam optimizer \citep{kingma2015adam} with default momentum parameters ($\beta_1 = 0.9$, $\beta_2 = 0.999$). To ensure a fair comparison, AM, MMFM$_{\rm van}$ and MMFM$_{\rm const}$ share the same network architecture and training hyperparameters, with the only difference being the input dimension. We used 100 function evaluations for the low-dimensional experiments (note the official MMFM repository used 1001 evaluations). For EB5, CITE5, and CITE50, we used 2000 function evaluations. The low-dimensional datasets were implemented on an M2 MacBook Air, the EB5, CITE5 and CITE50 experiments were run on a single NVIDIA H200 GPU.



\begin{table*}[h]
\centering
\small
\setlength{\tabcolsep}{5pt}
\caption{Main hyperparameters for training the conditional Gaussian sampler.}
\label{tab:params-conditional-sampler}
\resizebox{0.65\textwidth}{!}{
\begin{tabular}{lccccc}
\toprule
 & \textbf{GoM} & \textbf{LV} & \textbf{EB5} & \textbf{CITE5} & \textbf{CITE50}\\
 \midrule
Hidden dim      & 64        & 64     & 1024   & 512    & 1024   \\
Layers          & 1            & 1      & 4      & 2      & 2      \\
Batch size          & 111          & 50      & 256      & 256      & 256      \\
LR ($\times 10^{-1}$) & $1$ & $1$ & $0.1$ & $0.1$ & $0.1$ \\
Train steps     & 300  & 300 & 20\,000 & 10\,000 & 10\,000 \\
\end{tabular}
}
\end{table*}



\begin{table*}[h]
\centering
\small
\setlength{\tabcolsep}{5pt}
\caption{Collection of hyperparameters per dataset.}
\label{tab:params-acceleration}
\resizebox{0.65\textwidth}{!}{
\begin{tabular}{lccccc}
\toprule
 & \textbf{GoM} & \textbf{LV} & \textbf{EB5} & \textbf{CITE5} & \textbf{CITE50}\\
\midrule
\multicolumn{6}{l}{\textit{Reference process}} \\[2pt]

$\sigma^2_v$   & 50   & 50  & 0.005 & 0.01  & 1     \\
$\sqrt{\varepsilon}$  & 4     & 2   & 0.2   & 0.01 & 0.01  \\
\midrule
\multicolumn{6}{l}{\textit{Network and training parameters}} \\[2pt]
Architecture    & MLP    & MLP     & MLP & MLP & MLP    \\
Hidden dim      & 256    & 256     & 256    & 256    & 1024    \\
Layers          & 2      & 2       & 5     & 5   & 4      \\
Batch size      & 111    & 50      & 256     & 256      & 256      \\
LR $(\times 10^{-2})$    & 1 & 1  & 1    & 0.1 & 0.1 \\
Train steps     & 300   &300       & 20\,000 & 20\,000 & 20\,000 \\
\end{tabular}}
\end{table*}





\end{document}

